
\Data\ID{1003030201}
\Updated{2021-01-14 03:01:32}
\Level{A}
\SubArea{Lines and planes: intersecting, perpendicular, parallel}
\LANGen{
\Question{
Which of the following mathematical sentences corresponds to  the statement: Point \( C \) lies on the ray \( AB \)?}
{\( C \in\,\, \mapsto AB \)}{\( C \subset\,\, \mapsto AB \)}{\( C\,\,\cap\mapsto AB \)}{\( C\equiv\,\,\mapsto AB \)}{\( C \cong\,\,\mapsto AB \)}{}}
\LANGpl{
\Question{
Wskaż matematyczny zapis zdania: Punkt  \( C \) leży na półprostej \( AB \)?}
{\( C \in\,\, \mapsto AB \)}{\( C \subset\,\, \mapsto AB \)}{\( C\,\,\cap\mapsto AB \)}{\( C\equiv\,\,\mapsto AB \)}{\( C \cong\,\,\mapsto AB \)}{}}
\LANGcs{
\Question{
Jaký je symbolický zápis tvrzení "bod \( C \) leží na polopřímce \( AB \)"?}
{\( C \in\,\, \mapsto AB \)}{\( C \subset\,\, \mapsto AB \)}{\( C\,\,\cap\mapsto AB \)}{\( C\equiv\,\,\mapsto AB \)}{\( C \cong\,\,\mapsto AB \)}{}}
\LANGsk{
\Question{
Aký je symbolický zápis tvrdení "bod \( C \) leží na polopriame \( AB \)"?}
{\( C \in\,\, \mapsto AB \)}{\( C \subset\,\, \mapsto AB \)}{\( C\,\,\cap\mapsto AB \)}{\( C\equiv\,\,\mapsto AB \)}{\( C \cong\,\,\mapsto AB \)}{}}
\LANGes{
\Question{
¿Cuál es la notación simbólica de "punto \( C \) pertenece a la semirecta \( AB \)"?}
{\( C \in\,\, \mapsto AB \)}{\( C \subset\,\, \mapsto AB \)}{\( C\,\,\cap\mapsto AB \)}{\( C\equiv\,\,\mapsto AB \)}{\( C \cong\,\,\mapsto AB \)}{}}
\End

\Data\ID{1003030202}
\Updated{2021-01-14 03:01:02}
\Level{A}
\SubArea{Lines and planes: intersecting, perpendicular, parallel}
\LANGen{
\Question{
Which of the following mathematical sentences corresponds to  the statement: The segments \( AC \) and \( BD \) intersect at only one point \( C \)?}
{\( AC\cap BD=\{C\} \)}{\( AC\cap BD\subset\{C\} \)}{\( AC\cup BD=\{C\} \)}{\( AC\cup BD\subset\{C\} \)}{}{}}
\LANGpl{
\Question{
Wskaż matematyczny zapis zdania: Odcinki \( AC \) i \( BD \) mają tylko jeden punkt wspólny \( C \)?}
{\( AC\cap BD=\{C\} \)}{\( AC\cap BD\subset\{C\} \)}{\( AC\cup BD=\{C\} \)}{\( AC\cup BD\subset\{C\} \)}{}{}}
\LANGcs{
\Question{
Jaký je symbolický zápis tvrzení "úsečky \( AC \) a \( BD \) mají jediný společný bod  \( C \)"?}
{\( AC\cap BD=\{C\} \)}{\( AC\cap BD\subset\{C\} \)}{\( AC\cup BD=\{C\} \)}{\( AC\cup BD\subset\{C\} \)}{}{}}
\LANGsk{
\Question{
Aký je symbolický zápis tvrdenia "úsečky \( AC \) a \( BD \) majú jediný spoločný bod  \( C \)"?}
{\( AC\cap BD=\{C\} \)}{\( AC\cap BD\subset\{C\} \)}{\( AC\cup BD=\{C\} \)}{\( AC\cup BD\subset\{C\} \)}{}{}}
\LANGes{
\Question{
¿Cuál es la notación simbólica de  "segmentos \( AC \) y \( BD \) tienen en común un punto \( C \)"?}
{\( AC\cap BD=\{C\} \)}{\( AC\cap BD\subset\{C\} \)}{\( AC\cup BD=\{C\} \)}{\( AC\cup BD\subset\{C\} \)}{}{}}
\End

\Data\ID{1003030203}
\Updated{2021-01-14 03:01:57}
\Level{A}
\SubArea{Lines and planes: intersecting, perpendicular, parallel}
\LANGen{
\Question{
Suppose there are \( n \) points that lie on a line. Into how many not overlapping segments is the line divided by the points?}
{\( ( n-1 ) \) segments and two rays}{\( n \) segments}{\( (n+1) \)   segments}{\( n \) segments and two rays}{\( (n+1) \) segments and two rays}{}}
\LANGpl{
\Question{
Załóżmy, że \( n \) punktów leży na prostej.  Na ile niepokrywających się odcinków prosta jest podzielona przez punkty?}
{\( ( n-1 ) \) odcinków i dwie półproste}{\( n \) odcinków}{\( (n+1) \)   odcinków}{\( n \) odcinków i dwie półproste}{\( (n+1) \) odcinków i dwie półproste}{}}
\LANGcs{
\Question{
Na kolik vzájemně nepřekrývajících se částí rozdělí přímku \( n \) různých bodů, které na ní leží?}
{\( ( n-1 ) \) úseček a dvě polopřímky}{\( n \) úseček}{\( (n+1) \) úseček}{\( n \) úseček a dvě polopřímky}{\( (n+1) \) úseček a dvě polopřímky}{}}
\LANGsk{
\Question{
Na koľko vzájomne neprekrývajúcich sa častí rozdelia priamku \( n \) rôzne body, ktoré na nej ležia?}
{\( ( n-1 ) \) úsečiek a dve polpriamky}{\( n \) úsečiek}{\( (n+1) \) úsečiek}{\( n \) úsečiek a dve polpriamky}{\( (n+1) \) úsečiek a dve polpriamky}{}}
\LANGes{
\Question{
¿A cuántas partes disjuntas se divide una recta por \( n \) puntos distintos pertenecientes a dicha recta?}
{\( ( n-1 ) \) segmentos y dos semirectas}{\( n \) segmentos}{\( (n+1) \) segmentos}{\( n \) segmentos y dos semirectas}{\( (n+1) \) segmentos y dos semirectas}{}}
\End

\Data\ID{1003030204}
\Updated{2021-01-14 04:01:27}
\Level{A}
\SubArea{Lines and planes: intersecting, perpendicular, parallel}
\LANGen{
\Question{
What do we call two lines which lie on a single plane and have no point of intersection?}
{Distinct parallel lines}{Intersecting lines}{Identical lines}{Skew lines}{}{}}
\LANGpl{
\Question{
Jak nazywamy dwie proste na  płaszczyźnie, które  nie maja punktów wspólnych?}
{Proste równoległe, niepokrywające się}{Proste przecinające się}{Proste pokrywające się}{Proste skośne}{}{}}
\LANGcs{
\Question{
Dvě přímky ležící v jedné rovině, které nemají žádný společný bod, se nazývají:}
{různé rovnoběžné přímky}{různoběžné přímky}{totožné přímky}{mimoběžné přímky}{}{}}
\LANGsk{
\Question{
Dve priamky ležiace v jednej rovine, ktoré nemajú žiadny spoločný bod, sa nazývajú:}
{rôzne rovnobežné priamky}{rôznobežné priamky}{totožné priamky}{mimobežné priamky}{}{}}
\LANGes{
\Question{
Dos rectas pertenecientes al mismo plano que no tienen nigngún punto en común se llaman:}
{rectas paralelas distintas}{rectas secantes}{rectas coincidentes}{rectas cruzantes}{}{}}
\End

\Data\ID{1003030205}
\Updated{2021-01-14 04:01:23}
\Level{A}
\SubArea{Lines and planes: intersecting, perpendicular, parallel}
\LANGen{
\Question{
What do we call two lines which lie on a single plane and share one point of intersection?}
{Intersecting lines}{Distinct parallel lines}{Identical lines}{Skew lines}{}{}}
\LANGpl{
\Question{
Jak nazywamy proste na płaszczyźnie, które maja jeden punkt wspólny?}
{Proste przecinające się}{Proste równoległe, niepokrywające się}{Proste pokrywające się}{Proste skośne}{}{}}
\LANGcs{
\Question{
Dvě přímky ležící v jedné rovině, které mají právě jeden společný bod, se nazývají:}
{různoběžné přímky}{různé rovnoběžné přímky}{totožné přímky}{mimoběžné přímky}{}{}}
\LANGsk{
\Question{
Dve priamky ležiace v jednej rovine, ktoré majú práve jeden spoločný bod, sa nazývajú:}
{rôznobežné priamky}{rôzne rovnobežné priamky}{totožné priamky}{mimobežné priamky}{}{}}
\LANGes{
\Question{
Dos rectas que pertenecen al mismo plano y tienen exactamente un punto en común se llaman:}
{rectas secantes}{rectas paralelas distintas}{rectas coincidentes}{rectas cruzantes}{}{}}
\End

\Data\ID{1003030206}
\Updated{2021-01-14 04:01:10}
\Level{A}
\SubArea{Lines and planes: intersecting, perpendicular, parallel}
\LANGen{
\Question{
What do we call two lines which share all their points?}
{Identical lines}{Distinct parallel lines}{Skew lines}{Intersecting lines}{}{}}
\LANGpl{
\Question{
Jak nazywamy proste, które mają wszystkie punkty wspólne?}
{Proste pokrywające się}{Proste równoległe, niepokrywające się}{Proste skośne}{Proste przecinające się}{}{}}
\LANGcs{
\Question{
Dvě přímky, které mají společné všechny své body, se nazývají:}
{totožné přímky}{různé rovnoběžné přímky}{mimoběžné přímky}{různoběžné přímky}{}{}}
\LANGsk{
\Question{
Dve priamky, ktoré majú spoločné všetky svoje body, sa nazývajú:}
{totožné priamky}{rôzne rovnobežné priamky}{mimobežné priamky}{rôznobežné priamky}{}{}}
\LANGes{
\Question{
Dos rectas que comparten todos sus puntos se llaman:}
{rectas coincidentes}{rectas paralelas distintas}{rectas cruzantes}{rectas secantes}{}{}}
\End

\Data\ID{1003030207}
\Updated{2021-01-14 04:01:39}
\Level{A}
\SubArea{Lines and planes: intersecting, perpendicular, parallel}
\LANGen{
\Question{
There are two distinct parallel lines on one plane: \( p =\,\, \leftrightarrow KL \)  and \( q=\,\,\leftrightarrow MN \). What is the intersection of the half plane \( KLM \) and the half plane \( MNL \)?
(Please note that a hypothetical half plane \( XYZ \) is bounded by a line \( XY \) and contains a point \( Z \).)}
{A plane zone limited by the lines \( p \) and \( q \)}{The half plane \( KLM \)}{The convex angle \( LKM \)}{The quadrilateral \( KLMN \)}{The triangle \( MNL \)}{}}
\LANGpl{
\Question{
Na jednej płaszczyźnie istnieją dwie różne proste równoległe: \( p =\,\, \leftrightarrow KL \)  oraz  \( q=\,\,\leftrightarrow MN \). Przecięcie półpłaszczyzny \( KLM \) i \( MNL \) to?
(Należy pamiętać, że półpłaszczyzna \( XYZ \) jest ograniczona prostą \( XY \) i zawiera punkt  \( Z \).)}
{Obszar płaszczyzny ograniczony prostymi  \( p \) i \( q \)}{Półpłaszczyzna \( KLM \)}{Kąt wypukły \( LKM \)}{Czworokąt \( KLMN \)}{Trójkąt \( MNL \)}{}}
\LANGcs{
\Question{
V rovině jsou dány dvě různé rovnoběžné přímky \( p =\,\, \leftrightarrow KL \)  a \( q=\,\,\leftrightarrow MN \). Průnikem poloroviny \( KLM \) a poloroviny \( MNL \) je:
(Připomeňme, že polorovinou \( XYZ \) rozumíme polorovinu, která je ohraničená přímkou \( XY \) a obsahuje bod \( Z \).)}
{rovinný pás ohraničený přímkami \( p \) a \( q \)}{polorovina \( KLM \)}{konvexní úhel \( LKM \)}{čtyřúhelník \( KLMN \)}{trojúhelník \( MNL \)}{}}
\LANGsk{
\Question{
V rovine sú dané dve rôzne rovnobežné priamky \( p =\,\, \leftrightarrow KL \)  a \( q=\,\,\leftrightarrow MN \). Prienikom polroviny \( KLM \) a polroviny \( MNL \) je:
(Pripomeňme, že polrovinou \( XYZ \) rozumieme polrovinu, ktorá je ohraničená priamkou \( XY \) a obsahuje bod \( Z \).)}
{rovinný pás ohraničený priamkami \( p \) a \( q \)}{polrovina \( KLM \)}{konvexný úhol \( LKM \)}{štvoruholník \( KLMN \)}{trojuholník \( MNL \)}{}}
\LANGes{
\Question{
En el plano tenemos dos rectas paralelas distintas\( p =\,\, \leftrightarrow KL \)  y \( q=\,\,\leftrightarrow MN \). La intersección del semiplano \( KLM \) y el semiplano \( MNL \) es:
(Notamos que por el semiplano \( XYZ \) entendemos al semiplano acotado por la recta \( XY \) que contiene al punto \( Z \).)}
{cinta de plano acotada por las rectas \( p \) y \( q \)}{semiplano \( KLM \)}{ángulo convexo \( LKM \)}{cuadrilátero \( KLMN \)}{triángulo \( MNL \)}{}}
\End

\Data\ID{1003030208}
\Updated{2021-01-14 04:01:46}
\Level{A}
\SubArea{Lines and planes: intersecting, perpendicular, parallel}
\LANGen{
\Question{
What do we call a convex angle which is smaller than the right angle?}
{Acute angle}{Obtuse angle}{Straight angle}{Full angle}{}{}}
\LANGpl{
\Question{
Jak nazywamy kąt wypukły który jest mniejszy od kąta prostego?}
{Kąt ostry}{Kąt rozwarty}{Kąt półpełny}{Kat pełny}{}{}}
\LANGcs{
\Question{
Konvexní úhel, který je menší než pravý, se nazývá:}
{ostrý úhel}{tupý úhel}{přímý úhel}{plný úhel}{}{}}
\LANGsk{
\Question{
Konvexný uhol, ktorý je menší než pravý, sa nazýva:}
{ostrý uhol}{tupý uhol}{priamy uhol}{plný uhol}{}{}}
\LANGes{
\Question{
El ángulo convexo menor que el ángulo recto se llama:}
{ángulo agudo}{ángulo obtuso}{ángulo llano}{ángulo completo}{}{}}
\End

\Data\ID{1003030209}
\Updated{2021-01-14 04:01:38}
\Level{A}
\SubArea{Lines and planes: intersecting, perpendicular, parallel}
\LANGen{
\Question{
Let there be three points \( K \), \( L \), and \( M \) on a plane which do not lie on a single line. What is the intersection of the convex angles \( KLM \) and \( KML \)?
(Please note that a hypothetical angle \( XYZ \) is a part of a plane bounded by half lines \( YX \) and \( YZ \).)}
{The triangle \( KLM \)}{The segment \( KL \)}{The line \( p =\,\, \leftrightarrow KM \)}{The point \( K \)}{The axis of the convex angle \( KLM \)}{}}
\LANGpl{
\Question{
Dane są trzy punkty \( K \), \( L \) i \( M \) na płaszczyźnie,  punkty nie leżą na jednej prostej. Jakie jest przecięcie kątów wypukłych \( KLM \) i \( KML \)?
(Kąt  \( XYZ \) jest częścią płaszczyzny, która jest ograniczona przez półproste \( YX \) i \( YZ \).)}
{Trójkąt \( KLM \)}{Odcinek \( KL \)}{Prosta \( p =\,\, \leftrightarrow KM \)}{Punkt \( K \)}{Oś kąta wypukłego \( KLM \)}{}}
\LANGcs{
\Question{
V rovině jsou dány tři různé body \( K \), \( L \), \( M \) neležící v jedné přímce. Průnikem konvexních úhlů \( KLM \) a \( KML \) je:
(Připomeňme, že úhlem \( XYZ \) rozumíme část roviny, která je ohraničená polopřímkami \( YX \) a \( YZ \).)}
{trojúhelník \( KLM \)}{úsečka \( KL \)}{přímka \( p =\,\, \leftrightarrow KM \)}{bod \( K \)}{osa konvexního úhlu \( KLM \)}{}}
\LANGsk{
\Question{
V rovine sú dané tri rôzne body \( K \), \( L \), \( M \) , ktoré neležia v jednej priamke. Prienikom konvexných uhlov \( KLM \) a \( KML \) je:
(Pripomeňme, že uhlom \( XYZ \) rozumieme časť roviny, ktorá je ohraničená polpriamkami \( YX \) a \( YZ \).)}
{trojuholník \( KLM \)}{úsečka \( KL \)}{priamka \( p =\,\, \leftrightarrow KM \)}{bod \( K \)}{os konvexného uhla \( KLM \)}{}}
\LANGes{
\Question{
En el plano hay tres puntos distintos \( K \), \( L \), \( M \) que no pertenecen a una recta. Intersección del los ángulos conexos \( KLM \) y \( KML \) es:
(Připomeňme, že úhlem \( XYZ \) rozumíme část roviny, která je ohraničená polopřímkami \( YX \) a \( YZ \).)}
{triángulo \( KLM \)}{segmento \( KL \)}{recta \( p =\,\, \leftrightarrow KM \)}{punto \( K \)}{eje del ángulo convexo \( KLM \)}{}}
\End

\Data\ID{1003030301}
\Updated{2021-01-14 04:01:19}
\Level{A}
\SubArea{Lines and planes: intersecting, perpendicular, parallel}
\LANGen{
\Question{
Let \( p \) be a line and \( k \) be a circle with the centre \( S \) and the radius of \( 3\,\mathrm{cm} \). Given the distance of \( p \) and \( S \) is \( 4\,\mathrm{cm} \), the line \( p \) is:}
{an outer line of the circle \( k \)}{a secant of the circle \( k \)}{a tangent of the circle \( k \)}{a chord of the circle \( k \)}{}{}}
\LANGpl{
\Question{
Dana jest prosta \( p \) i okrąg \( k \) o środku \( S \) i promieniu równym \( 3\,\mathrm{cm} \). Odległość między  \( p \) i \( S \) wynosi \( 4\,\mathrm{cm} \), prosta \( p \) to:}
{prosta leżąca poza okręgiem \( k \)}{sieczna okręgu \( k \)}{styczna do okręgu \( k \)}{cięciwa okręgu \( k \)}{}{}}
\LANGcs{
\Question{
Je dána kružnice \( k( S;3\,\mathrm{cm})\) a přímka \( p \) tak, že \( |Sp|=4\,\mathrm{cm}\). Přímka  \( p \) je:}
{vnější přímka kružnice \( k \)}{sečna kružnice \( k \)}{tečna kružnice \( k \)}{tětiva kružnice \( k \)}{}{}}
\LANGsk{
\Question{
Je daná kružnica \( k( S;3\,\mathrm{cm})\) a priamka \( p \) tak, že \( |Sp|=4\,\mathrm{cm}\). Priamka  \( p \) je:}
{vonkajšia priamka kružnice \( k \)}{sečnica kružnice \( k \)}{dotyčnica kružnice \( k \)}{tetiva kružnice \( k \)}{}{}}
\LANGes{
\Question{
Sea la circunferencia \( k( S;3\,\mathrm{cm})\) y la recta \( p \) tales que \( |Sp|=4\,\mathrm{cm}\). La recta  \( p \) es:}
{recta externa de la circunferencia \( k \)}{recta secante a la circunferencia \( k \)}{recta tangente a la circunferencia \( k \)}{cuerda de la circunferencia \( k \)}{}{}}
\End

\Data\ID{1003030302}
\Updated{2021-01-14 04:01:57}
\Level{A}
\SubArea{Lines and planes: intersecting, perpendicular, parallel}
\LANGen{
\Question{
Let  \( p \) be a line and \( k \) be a circle with the centre \( S \) and the radius of \( 3\,\mathrm{cm} \). Given the distance of \( p \) and \( S \) is \( 2.8\,\mathrm{cm} \), the line \( p \) is:}
{a secant of the circle \( k \)}{a tangent of the circle \( k \)}{an outer line of the circle \( k \)}{a chord of the circle \( k \)}{}{}}
\LANGpl{
\Question{
Dano prostą   \( p \) i okrąg \( k \) o środku \( S \) i promieniu o długości \( 3\,\mathrm{cm} \). Odległość pomiędzy  \( p \) i \( S \) jest równa \( 2{,}8\,\mathrm{cm} \), prosta \( p \) to:}
{sieczna okręgu \( k \)}{styczna do okręgu \( k \)}{prosta poza okręgiem \( k \)}{cięciwa okręgu \( k \)}{}{}}
\LANGcs{
\Question{
Je dána kružnice \( k( S;3\,\mathrm{cm})\) a přímka  \( p \) tak, že \( |Sp|=2,8\,\mathrm{cm}\). Přímka \( p \) je:}
{sečna kružnice \( k \)}{tečna kružnice \( k \)}{vnější přímka kružnice \( k \)}{tětiva kružnice \( k \)}{}{}}
\LANGsk{
\Question{
Je daná kružnica \( k( S;3\,\mathrm{cm})\) a priamka  \( p \) tak, že \( |Sp|=2,8\,\mathrm{cm}\). Priamka \( p \) je:}
{sečnica kružnice \( k \)}{dotyčnica kružnice \( k \)}{vonkajšia priamka kružnice \( k \)}{tetiva kružnice \( k \)}{}{}}
\LANGes{
\Question{
Sea circunferencia \( k( S;3\,\mathrm{cm})\) y recta \( p \) tales que \( |Sp|=2,8\,\mathrm{cm}\). La recta \( p \) es:}
{secante a la circunferencia \( k \)}{tangente a la circunferencia \( k \)}{exterior a la circunferencia \( k \)}{cuerda de la circunferencia \( k \)}{}{}}
\End

\Data\ID{1003030303}
\Updated{2021-01-14 04:01:14}
\Level{A}
\SubArea{Lines and planes: intersecting, perpendicular, parallel}
\LANGen{
\Question{
Let \( k_1 \) and \( k_2 \) be circles with the centres \( S_1\), \( S_2 \),  and the radii of lengths \( 5\,\mathrm{cm} \) and \( 1\,\mathrm{cm} \) consecutively. Given \( |S_1 S_2|=3\,\mathrm{cm} \), what is the mutual position of these circles?}
{The circle \( k_2 \)  lies inside the circle  \( k_1 \).}{The circle  \( k_1 \) lies inside the circle  \( k_2 \).}{The circles \( k_1 \)   and \( k_2 \)   intersect.}{The circles  \( k_1 \)  and  \( k_2 \)  are externally tangent.}{The circles \( k_1 \)   a \( k_2 \)   are internally tangent.}{}}
\LANGpl{
\Question{
Niech \( k_1 \) i \( k_2 \) to okręgi o środkach \( S_1\), \( S_2 \),  długość ich promieni wynosi \( 5\,\mathrm{cm} \) i \( 1\,\mathrm{cm} \). Wiadomo, że  \( |S_1 S_2|=3\,\mathrm{cm} \), określ wzajemne położenie okręgów.}
{Okrąg \( k_2 \)  leży wewnątrz okręgu \( k_1 \).}{Okrąg  \( k_1 \)   leży wewnątrz okręgu  \( k_2 \).}{Okręgi \( k_1 \)   i \( k_2 \)   to okręgi przecinające się.}{Okręgi \( k_1 \)  i  \( k_2 \)  to okręgi zewnętrznie styczne.}{Okręgi \( k_1 \)   a \( k_2 \)   to okręgi wewnętrznie styczne.}{}}
\LANGcs{
\Question{
Jsou dány kružnice \( k_1( S_1;5\,\mathrm{cm})\), \( k_2( S_2;1\,\mathrm{cm})\), \( |S_1 S_2|=3\,\mathrm{cm} \). Jaká je vzájemná poloha těchto kružnic?}
{Kružnice \( k_2 \)  leží uvnitř kružnice  \( k_1 \).}{Kružnice  \( k_1 \) leží uvnitř kružnice  \( k_2 \).}{Kružnice \( k_1 \)   a \( k_2 \)  se protínají.}{Kružnice  \( k_1 \)  a  \( k_2 \)  mají vnější dotyk.}{Kružnice \( k_1 \)   a \( k_2 \)   mají vnitřní dotyk.}{}}
\LANGsk{
\Question{
Sú dané kružnice \( k_1( S_1;5\,\mathrm{cm})\), \( k_2( S_2;1\,\mathrm{cm})\), \( |S_1 S_2|=3\,\mathrm{cm} \). Aká je vzájomná poloha týchto kružníc?}
{Kružnica \( k_2 \)  leží vnútri kružnice  \( k_1 \).}{Kružnica  \( k_1 \) leží vnútri kružnice  \( k_2 \).}{Kružnice \( k_1 \)   a \( k_2 \)  sa pretínajú.}{Kružnice  \( k_1 \)  a  \( k_2 \)  majú vonkajší dotyk.}{Kružnice \( k_1 \)   a \( k_2 \)   majú vnútorný dotyk.}{}}
\LANGes{
\Question{
Sean las ircunferencias \( k_1( S_1;5\,\mathrm{cm})\), \( k_2( S_2;1\,\mathrm{cm})\), \( |S_1 S_2|=3\,\mathrm{cm} \). ¿Cuál es su posición?}
{La circunferencia \( k_2 \) está dentro de la circunferencia \( k_1 \).}{La circunferencia \( k_1 \) está dentro de la circunferencia \( k_2 \).}{Las circunferencias \( k_1 \)  y \( k_2 \)  se secan.}{Las circunferencias \( k_1 \)  y  \( k_2 \)  son tangentes exteriores.}{Las circunferencias \( k_1 \)  y \( k_2 \)  son tangentes interiores.}{}}
\End

\Data\ID{1003030304}
\Updated{2021-01-14 04:01:19}
\Level{A}
\SubArea{Lines and planes: intersecting, perpendicular, parallel}
\LANGen{
\Question{
Let \( k_1 \) and \( k_2 \) be circles with the centres \( S_1 \), \( S_2 \),  and the radii of lengths \( 5\,\mathrm{cm} \) and \( 1\,\mathrm{cm} \) consecutively. Given \( |S_1 S_2 |=7\,\mathrm{cm} \), what is the mutual position of these circles?}
{The circle \( k_2 \)  lies outside the circle  \( k_1 \).}{The circle \( k_2 \)  lies inside the circle  \( k_1 \).}{The circles  \( k_1 \)  and  \( k_2 \)  intersect.}{The circles \( k_1 \)   and \( k_2 \)   are externally tangent.}{The circles \( k_1 \)   and  \( k_2 \)  are internally tangent.}{}}
\LANGpl{
\Question{
Niech  \( k_1 \) i \( k_2 \) to okręgi o środkach  \( S_1 \), \( S_2 \),  długość ich promieni jest równa \( 5\,\mathrm{cm} \) i \( 1\,\mathrm{cm} \). Wiadomo, że  \( |S_1 S_2 |=7\,\mathrm{cm} \), określ wzajemne położenie okręgów.}
{Okrąg \( k_2 \)   leży na zewnątrz okręgu  \( k_1 \).}{Okrąg \( k_2 \)   leży wewnątrz okręgu \( k_1 \).}{Okręgi \( k_1 \)  i  \( k_2 \)  to okręgi przecinające się.}{Okręgi \( k_1 \)   i \( k_2 \)  to okręgi zewnętrznie styczne.}{Okręgi \( k_1 \)   i  \( k_2 \)  to okręgi wewnętrznie styczne.}{}}
\LANGcs{
\Question{
Jsou dány kružnice \( k_1( S_1;5\,\mathrm{cm})\), \( k_2( S_2;1\,\mathrm{cm})\), \( |S_1 S_2 |=7\,\mathrm{cm} \). Jaká je vzájemná poloha těchto kružnic?}
{Kružnice \( k_2 \)  leží vně kružnice  \( k_1 \).}{Kružnice \( k_2 \)  leží uvnitř kružnice  \( k_1 \).}{Kružnice  \( k_1 \)  a  \( k_2 \)  se protínají.}{Kružnice \( k_1 \)  a \( k_2 \)  mají vnější dotyk.}{Kružnice \( k_1 \)   a  \( k_2 \)  mají vnitřní dotyk.}{}}
\LANGsk{
\Question{
Sú dané kružnice \( k_1( S_1;5\,\mathrm{cm})\), \( k_2( S_2;1\,\mathrm{cm})\), \( |S_1 S_2 |=7\,\mathrm{cm} \). Aká je vzájomná poloha týchto kružníc?}
{Kružnica \( k_2 \)  leží mimo kružnice  \( k_1 \).}{Kružnica \( k_2 \)  leží vnútri kružnice  \( k_1 \).}{Kružnice  \( k_1 \)  a  \( k_2 \)  sa pretínajú.}{Kružnice \( k_1 \)  a \( k_2 \)  majú vonkajší dotyk.}{Kružnice \( k_1 \)   a  \( k_2 \)  majú vnútorný dotyk.}{}}
\LANGes{
\Question{
Sea dos circunferencias \( k_1( S_1;5\,\mathrm{cm})\), \( k_2( S_2;1\,\mathrm{cm})\), \( |S_1 S_2 |=7\,\mathrm{cm} \). ¿Cuál es su posición?}
{La circunferencia \( k_2 \)  es exterior a la circunferencia \( k_1 \).}{La circunferencia \( k_2 \) es interior a la circunferencia \( k_1 \).}{Las circunferencias \( k_1 \)  y  \( k_2 \)  se secan.}{Las circunferencias \( k_1 \)  y \( k_2 \)  son secantes exteriores.}{Las circunferencias \( k_1 \)  y  \( k_2 \)  son secantes interiores.}{}}
\End

\Data\ID{1003030305}
\Updated{2021-01-14 04:01:54}
\Level{A}
\SubArea{Lines and planes: intersecting, perpendicular, parallel}
\LANGen{
\Question{
Let \( k_1 \) and \( k_2 \) be circles with the centres \( S_1 \), \( S_2 \),  and the radii of lengths \( 5\,\mathrm{cm} \) and \( 1\,\mathrm{cm} \) consecutively. Given \( |S_1S_2|=6\,\mathrm{cm}\), what is the mutual position of these circles?}
{The circles \( k_1 \)    and \( k_2 \)   are externally tangent.}{The circles  \( k_1 \)  and  \( k_2 \)  are internally tangent.}{The circle \( k_1 \)  lies inside the circle  \( k_2 \).}{The circle \( k_2 \)  lies inside the circle  \( k_1 \).}{}{}}
\LANGpl{
\Question{
Niech \( k_1 \) to  \( k_2 \) to okręgi o środkach \( S_1 \), \( S_2 \),  a długość ich promieni wynosi \( 5\,\mathrm{cm} \) i \( 1\,\mathrm{cm} \). Wiadomo, że  \( |S_1S_2|=6\,\mathrm{cm}\), określ wzajemne położenie okręgów.}
{Okręgi \( k_1 \)    i  \( k_2 \)   to okręgi zewnętrznie styczne.}{Okręgi  \( k_1 \)  i  \( k_2 \)  to okręgi wewnętrznie styczne.}{Okrąg \( k_1 \)   leży wewnątrz okręgu   \( k_2 \).}{Okrąg  \( k_2 \)   leży wewnątrz okręgu   \( k_1 \).}{}{}}
\LANGcs{
\Question{
Jsou dány kružnice \( k_1( S_1;5\,\mathrm{cm})\), \( k_2( S_2;1\,\mathrm{cm})\),  \( |S_1S_2|=6\,\mathrm{cm}\). Jaká je vzájemná poloha těchto kružnic?}
{Kružnice \( k_1 \)   a  \( k_2 \)   mají vnější dotyk.}{Kružnice  \( k_1 \)  a  \( k_2 \)  mají vnitřní dotyk.}{Kružnice  \( k_1 \)  leží uvnitř kružnice  \( k_2 \).}{Kružnice \( k_2 \)  leží uvnitř kružnice  \( k_1 \).}{}{}}
\LANGsk{
\Question{
Sú dané kružnice \( k_1( S_1;5\,\mathrm{cm})\), \( k_2( S_2;1\,\mathrm{cm})\),  \( |S_1S_2|=6\,\mathrm{cm}\). Aká je vzájomná poloha týchto kružníc?}
{Kružnice \( k_1 \)   a  \( k_2 \)   majú vonkajší dotyk.}{Kružnice  \( k_1 \)  a  \( k_2 \)  majú vnútorný dotyk.}{Kružnica  \( k_1 \)  leží vnútri kružnice  \( k_2 \).}{Kružnica \( k_2 \)  leží vnútri kružnice  \( k_1 \).}{}{}}
\LANGes{
\Question{
Sean las circunferencias \( k_1( S_1;5\,\mathrm{cm})\), \( k_2( S_2;1\,\mathrm{cm})\),  \( |S_1S_2|=6\,\mathrm{cm}\). ¿Cuál es su posición?}
{Las circinferencias \( k_1 \) y \( k_2 \) son tangentes exteriores.}{Las circinferencias \( k_1 \) y \( k_2 \) son tangentes interiores.}{La circunferencia \( k_1 \) está dentro de la circunferencia \( k_2 \).}{La circunferencia \( k_2 \) está dentro de la circunferencia \( k_1\).}{}{}}
\End

\Data\ID{1003030306}
\Updated{2021-01-14 04:01:37}
\Level{A}
\SubArea{Lines and planes: intersecting, perpendicular, parallel}
\LANGen{
\Question{
Let \( k_1 \) and \( k_2 \) be circles with the centres \( S_1 \), \( S_2 \),  and the radii of lengths \( 5\,\mathrm{cm} \) and \( 1\,\mathrm{cm} \) consecutively. Given \( |S_1S_2|=4\,\mathrm{cm}\), what is the mutual position of these circles?}
{The circles \( k_1 \)  and \( k_2 \) are internally tangent.}{The circles   \( k_1 \)   and    \( k_2 \)  are externally tangent.}{The circle  \( k_1 \)  lies inside the circle   \( k_2 \).}{The circle  \( k_2 \)  lies inside the circle   \( k_1 \).}{}{}}
\LANGpl{
\Question{
Niech  \( k_1 \) i \( k_2 \) to okręgi o środkach \( S_1 \), \( S_2 \), długość ich promieni jest równa \( 5\,\mathrm{cm} \) i \( 1\,\mathrm{cm} \). Wiadomo, że \( |S_1S_2|=4\,\mathrm{cm}\), określ wzajemne położenie okręgów.}
{Okręgi \( k_1 \)  i \( k_2 \) to okręgi wewnętrznie styczne.}{Okręgi   \( k_1 \)   i   \( k_2 \)  to okręgi zewnętrznie styczne.}{Okrąg  \( k_1 \)   leży wewnątrz okręgu  \( k_2 \).}{Okrąg \( k_2 \)   leży wewnątrz okręgu  \( k_1 \).}{}{}}
\LANGcs{
\Question{
Jsou dány kružnice \( k_1( S_1;5\,\mathrm{cm})\), \( k_2( S_2;1\,\mathrm{cm})\), \( |S_1S_2|=4\,\mathrm{cm}\). Jaká je vzájemná poloha těchto kružnic?}
{Kružnice \( k_1 \)  a  \( k_2 \) mají vnitřní dotyk.}{Kružnice   \( k_1 \)   a   \( k_2 \)  mají vnější dotyk.}{Kružnice  \( k_1 \)  leží uvnitř kružnice   \( k_2 \).}{Kružnice  \( k_2 \)  leží uvnitř kružnice   \( k_1 \).}{}{}}
\LANGsk{
\Question{
Sú dané kružnice \( k_1( S_1;5\,\mathrm{cm})\), \( k_2( S_2;1\,\mathrm{cm})\), \( |S_1S_2|=4\,\mathrm{cm}\). Aká je vzájomná poloha týchto kružníc?}
{Kružnice \( k_1 \)  a  \( k_2 \) majú vnútorný dotyk.}{Kružnice   \( k_1 \)   a   \( k_2 \)  majú vonkajší dotyk.}{Kružnica  \( k_1 \)  leží vnútri kružnice   \( k_2 \).}{Kružnica  \( k_2 \)  leží vnútri kružnice   \( k_1 \).}{}{}}
\LANGes{
\Question{
Sean las circunferencias \( k_1( S_1;5\,\mathrm{cm})\), \( k_2( S_2;1\,\mathrm{cm})\), \( |S_1S_2|=4\,\mathrm{cm}\). ¿Cuál es su posición?}
{Las circinferencias \( k_1 \) y \( k_2 \) son tangentes interiores.}{Las circinferencias \( k_1 \) y \( k_2 \) son tangentes exteriores.}{La circunferencia \( k_1 \) está dentro de la circunferencia \( k_2 \).}{La circunferencia \( k_2\) está dentro de la circunferencia \( k_1 \).}{}{}}
\End

\Data\ID{1003030307}
\Updated{2021-01-14 04:01:55}
\Level{A}
\SubArea{Lines and planes: intersecting, perpendicular, parallel}
\LANGen{
\Question{
Let \( k_1 \) and \( k_2 \) be circles with the centres \( S_1 \), \( S_2 \),  and the radii of lengths \( 5\,\mathrm{cm} \) and \( 1\,\mathrm{cm} \) consecutively. Given \( |S_1S_2|=5\,\mathrm{cm} \), what is the mutual position of these circles?}
{The circles  \( k_1 \)  and \( k_2 \)  intersect.}{The circles  \( k_1 \)  and \( k_2 \)   are internally tangent.}{The circles \( k_1 \)   and \( k_2 \)   are externally tangent.}{The circle \( k_1 \)  lies inside the circle  \( k_2 \).}{The circle \( k_2 \)  lies inside the circle  \( k_1 \).}{}}
\LANGpl{
\Question{
Niech  \( k_1 \) i \( k_2 \) to okręgi o środkach \( S_1 \), \( S_2 \),  a długość promieni okręgów jest równa  \( 5\,\mathrm{cm} \) i \( 1\,\mathrm{cm} \). Wiadomo, że \( |S_1S_2|=5\,\mathrm{cm} \), określ wzajemne położenie okręgów.}
{Okręgi  \( k_1 \)  i \( k_2 \)  to okręgi przecinające się.}{Okręgi  \( k_1 \)  i \( k_2 \)   to okręgi wewnętrznie styczne.}{Okręgi \( k_1 \)   i \( k_2 \)   to okręgi zewnętrznie styczne.}{Okrąg \( k_1 \)   leży wewnątrz okręgu  \( k_2 \).}{Okrąg \( k_2 \)   leży wewnątrz okręgu   \( k_1 \).}{}}
\LANGcs{
\Question{
Jsou dány kružnice \( k_1( S_1;5\,\mathrm{cm})\), \( k_2( S_2;1\,\mathrm{cm})\), \( |S_1S_2|=5\,\mathrm{cm} \). Jaká je vzájemná poloha těchto kružnic?}
{Kružnice  \( k_1 \)  a \( k_2 \)  se protínají.}{Kružnice  \( k_1 \)  a \( k_2 \)   mají vnitřní dotyk.}{Kružnice  \( k_1 \)   a  \( k_2 \)   mají vnější dotyk.}{Kružnice \( k_1 \)  leží uvnitř kružnice  \( k_2 \).}{Kružnice \( k_2 \)  leží uvnitř kružnice  \( k_1 \).}{}}
\LANGsk{
\Question{
Sú dané kružnice \( k_1( S_1;5\,\mathrm{cm})\), \( k_2( S_2;1\,\mathrm{cm})\), \( |S_1S_2|=5\,\mathrm{cm} \). Jaká je vzájomná poloha týchto kružníc?}
{Kružnice  \( k_1 \)  a \( k_2 \)  sa pretínajú.}{Kružnice  \( k_1 \)  a \( k_2 \)   majú vnútorný dotyk.}{Kružnice  \( k_1 \)   a  \( k_2 \)   majú vonkajší dotyk.}{Kružnica \( k_1 \)  leží vnútri kružnice  \( k_2 \).}{Kružnica \( k_2 \)  leží vnútri kružnice  \( k_1 \).}{}}
\LANGes{
\Question{
Sean las circunferencias \( k_1( S_1;5\,\mathrm{cm})\), \( k_2( S_2;1\,\mathrm{cm})\), \( |S_1S_2|=5\,\mathrm{cm} \). ¿Cuál es su posición?}
{Kružnice  \( k_1 \)  a \( k_2 \)  se protínají.}{Las circinferencias \( k_1 \) y \( k_2 \) son tangentes interiores.}{Las circinferencias \( k_1 \) y \( k_2 \) son tangentes exteriores.}{La circunferencia \( k_1 \) está dentro de la circunferencia \( k_2 \).}{La circunferencia \( k_2 \) está dentro de la circunferencia \( k_1 \).}{}}
\End

\Data\ID{1003030308}
\Updated{2021-01-14 04:01:47}
\Level{A}
\SubArea{Lines and planes: intersecting, perpendicular, parallel}
\LANGen{
\Question{
Let \( k_1 \) and \( k_2 \) be circles with the centres \( S_1 \), \( S_2 \),  and the radii of lengths \( 5\,\mathrm{cm} \) and \( 1\,\mathrm{cm} \) consecutively. Given \(S_1=S_2\), what is the mutual position of these circles?}
{The circles \( k_1 \)   and   \( k_2 \)   are concentric.}{The circles   \( k_1 \)  and  \( k_2 \)   intersect.}{The circle  \( k_1 \)   lies inside the circle  \( k_2 \).}{The circle   \( k_2 \) lies outside the circle   \( k_1 \).}{The circles  \( k_1 \)   and   \( k_2 \)  are internally tangent.}{}}
\LANGpl{
\Question{
Niech \( k_1 \) i \( k_2 \) to okręgi o środkach \( S_1 \), \( S_2 \),  długość promieni jest równa \( 5\,\mathrm{cm} \) i \( 1\,\mathrm{cm} \). Wiadomo, że  \(S_1=S_2\), określ wzajemne położenie tych okręgów?}
{Okręgi \( k_1 \)   i  \( k_2 \)   to okręgi współśrodkowe.}{Okręgi   \( k_1 \)  i  \( k_2 \)   to okręgi przecinające się.}{Okrąg  \( k_1 \)   leży wewnątrz okręgu  \( k_2 \).}{Okrąg   \( k_2 \) leży na zewnątrz okręgu   \( k_1 \).}{Okręgi \( k_1 \)   i   \( k_2 \)  są wewnętrznie styczne.}{}}
\LANGcs{
\Question{
Jsou dány kružnice \( k_1( S_1;5\,\mathrm{cm})\), \( k_2( S_2;1\,\mathrm{cm})\), \(S_1=S_2\). Jaká je vzájemná poloha těchto kružnic?}
{Kružnice \( k_1 \)   a   \( k_2 \)   jsou soustředné.}{Kružnice   \( k_1 \)  a  \( k_2 \)   se protínají.}{Kružnice  \( k_1 \)   leží uvnitř kružnice  \( k_2 \).}{Kružnice   \( k_2 \) leží vně kružnice   \( k_1 \).}{Kružnice  \( k_1 \)   a   \( k_2 \)  mají vnitřní dotyk.}{}}
\LANGsk{
\Question{
Sú dané kružnice \( k_1( S_1;5\,\mathrm{cm})\), \( k_2( S_2;1\,\mathrm{cm})\), \(S_1=S_2\). Aká je vzájomná poloha týchto kružníc?}
{Kružnice \( k_1 \)   a   \( k_2 \)   sú sústredné.}{Kružnice   \( k_1 \)  a  \( k_2 \)   sa pretínajú.}{Kružnica  \( k_1 \)   leží vnútri kružnice  \( k_2 \).}{Kružnica   \( k_2 \) leží mimo kružnicu   \( k_1 \).}{Kružnice  \( k_1 \)   a   \( k_2 \)  majú vnútorný dotyk.}{}}
\LANGes{
\Question{
Sean las circunferencias \( k_1( S_1;5\,\mathrm{cm})\), \( k_2( S_2;1\,\mathrm{cm})\), \(S_1=S_2\). ¿Cuál es su posición?}
{Las circunferencias \( k_1 \)  y  \( k_2 \) son concéntricas.}{Las circunferencias \( k_1 \)  y \( k_2 \)  se secan.}{La circunferencia \( k_1 \) está dentro de la circunferencia \( k_2 \).}{La circunferencia \( k_2\) está fuera de la circunferencia \( k_1\).}{Las circinferencias \( k_1 \) y \( k_2 \) son tangentes interiores.}{}}
\End

\Data\ID{1103030210}
\Updated{2021-01-15 09:01:28}
\Level{A}
\SubArea{Lines and planes: intersecting, perpendicular, parallel}
\IMG{1103030210.pdf}
\LANGen{
\Question{
Let there be two distinct parallel lines \( a \) and \( b \). What do we call the pair of in the picture  indicated angles \( \alpha \) and \( \beta \) which are defined by the transversal \( p \)?}
{Alternate angles}{Corresponding angles}{Adjacent angles}{Vertical angles}{}{}}
\LANGpl{
\Question{
Dane są dwie różne proste równoległe \( a \) i \( b \).  Jak nazywamy parę kątów  \( \alpha \) i  \( \beta \) zaznaczonych na rysunku, kąty  są określone linią poprzeczną  \( p \)?}
{Kąty naprzemianległe}{Kąty odpowiadające}{Kąty przyległe}{Kąty wierzchołkowe}{}{}}
\LANGcs{
\Question{
Jsou dány dvě různé rovnoběžné přímky \( a \), \( b \). Dvojice vyznačených úhlů \( \alpha \), \( \beta \) na obrázku, které jsou vyťaty příčkou \( p \) přímek  \( a \),  \( b \), se nazývají:}
{úhly střídavé}{úhly souhlasné}{úhly vedlejší}{úhly vrcholové}{}{}}
\LANGsk{
\Question{
Sú dané dve rôzne rovnobežné priamky \( a \), \( b \). Dvojice vyznačených uhlov \( \alpha \), \( \beta \) na obrázku, ktoré sú vyťaté priečkou \( p \) priamok  \( a \),  \( b \), sa nazývajú:}
{uhly striedavé}{uhly súhlasné}{uhly vedľajšie}{uhly vrcholové}{}{}}
\LANGes{
\Question{
Sean dos rectas paralelas \( a \), \( b \). La pareja de ángulos \( \alpha \), \( \beta \) marcados en el dibujo, definidos por la transversal \( p \) que corta a \( a \) y \( b \), se llama}
{ángulos alternos}{ángulos correspondientes}{ángulos adyacentes}{ángulos opuestos por el vértice}{}{}}
\End

\Data\ID{1003059501}
\Updated{2021-01-14 05:01:52}
\Level{B}
\SubArea{Lines and planes: intersecting, perpendicular, parallel}
\LANGen{
\Question{
Let \( \alpha \) and \( \beta \) be planes sharing three different points \( A \), \( B \), and \( C \). The points do not lie on one line. What is the mutual position of these two planes?}
{identical planes}{parallel planes}{intersecting planes}{}{}{}}
\LANGpl{
\Question{
Płaszczyzny \( \alpha \) i  \( \beta \) mają trzy różne wspólne punkty  \( A \), \( B \), i  \( C \). Punkty nie leżą na jednej prostej. Jakie jest wzajemne położenie tych dwóch płaszczyzn?}
{płaszczyzny identyczne}{płaszczyzny równoległe}{płaszczyzny przecinające się}{}{}{}}
\LANGcs{
\Question{
Jsou dány dvě roviny  \( \alpha \) a \( \beta \), které mají společné tři různé body \( A \), \( B \), \( C \) neležící v jedné přímce. Jaká je vzájemná poloha těchto rovin?}
{totožné roviny}{různé rovnoběžné roviny}{různoběžné roviny}{}{}{}}
\LANGsk{
\Question{
Sú dané dve roviny  \( \alpha \) a \( \beta \), ktoré majú spoločné tri rôzne body \( A \), \( B \), \( C \) neležiace na jednej priamke. Jaká je vzájomná poloha týchto rovín?}
{totožné roviny}{rôzne rovnobežné roviny}{rôznobežné roviny}{}{}{}}
\LANGes{
\Question{
Sean dos planos \( \alpha \) y \( \beta \), que tienen en común tres puntos distintos \( A \), \( B \), \( C \) no pertenecientes a una recta. ¿Cuál es la posición de los planos?}
{planos coincidentes}{planos paralelos}{planos secantes}{}{}{}}
\End

\Data\ID{1103059401}
\Updated{2021-01-15 10:01:22}
\Level{B}
\SubArea{Lines and planes: intersecting, perpendicular, parallel}
\IMG{1103059401.pdf}
\LANGen{
\Question{
Let \( ABCDEFGH \) be a cube and \( M \), \( N \), and \( P \) the midpoints of the edges \( BF \), \( EF \), and \( CD \), as seen in the picture. What is the mutual position of the lines \( DM \) and \( NP \)?}
{skew lines}{distinct parallel lines}{intersecting lines}{}{}{}}
\LANGpl{
\Question{
Dany jest sześcian \( ABCDEFGH \),  \( M \), \( N \), i \( P \) to środki krawędzi \( BF \), \( EF \), i \( CD \), zobacz rysunek. Jakie jest wzajemne położenie prostych \( DM \) i \( NP \)?}
{proste skośne}{proste równoległe, niepokrywające się}{proste przecinające się}{}{}{}}
\LANGcs{
\Question{
Je dána krychle \( ABCDEFGH \) a body \( M \), \( N \), \( P \), které jsou po řadě středy hran \( BF \), \( EF \) a \( CD \). Jaká je vzájemná poloha přímek \( DM \) a \( NP \)?}
{mimoběžné přímky}{různé rovnoběžné přímky}{různoběžné přímky}{}{}{}}
\LANGsk{
\Question{
Je daná kocka \( ABCDEFGH \) a body \( M \), \( N \), \( P \), ktoré sú po rade stredy hrán \( BF \), \( EF \) a \( CD \). Aká je vzájomná poloha priamok \( DM \) a \( NP \)?}
{mimobežné priamky}{rôzne rovnobežné priamky}{rôznobežné priamky}{}{}{}}
\LANGes{
\Question{
Sea un cubo \( ABCDEFGH \) y puntos \( M \), \( N \), \( P \), que son los medios de los bordes \( BF \), \( EF \) y \( CD \) respectamente. ¿Cuál es la pocisión de las rectas \( DM \) y \( NP \)?}
{rectas cruzantes}{rectas paralelas}{rectas secantes}{}{}{}}
\End

\Data\ID{1103059402}
\Updated{2021-01-15 10:01:39}
\Level{B}
\SubArea{Lines and planes: intersecting, perpendicular, parallel}
\IMG{1103059402.pdf}
\LANGen{
\Question{
Let \( ABCDEFGH \) be a cube and \( L \), \( M \) the midpoints of the edges \( AB \) and \( BF \). See the picture. What is the mutual position of lines \( DL \) and \( GM \)?}
{intersecting lines}{distinct parallel lines}{skew lines}{}{}{}}
\LANGpl{
\Question{
Dany jest sześcian \( ABCDEFGH \), \( L \), \( M \) to środki krawędzi  \( AB \) i \( BF \). Zobacz rysunek. Jakie jest wzajemne położenie prostych \( DL \) i \( GM \)?}
{proste przecinające się}{proste równoległe, niepokrywające się}{proste skośne}{}{}{}}
\LANGcs{
\Question{
Je dána krychle \( ABCDEFGH \) a body \( L \), \( M \), které jsou po řadě středy hran \( AB \) a \( BF \). Jaká je vzájemná poloha přímek \( DL \) a \( GM \)?}
{různoběžné přímky}{různé rovnoběžné přímky}{mimoběžné přímky}{}{}{}}
\LANGsk{
\Question{
Je daná kocka \( ABCDEFGH \) a body \( L \), \( M \), ktoré sú po rade stredy hrán \( AB \) a \( BF \). Aká je vzájomná poloha priamok \( DL \) a \( GM \)?}
{rôznobežné priamky}{rôzne rovnobežné priamky}{mimobežné priamky}{}{}{}}
\LANGes{
\Question{
Sea un cubo \( ABCDEFGH \) y los puntos \( L \), \( M \), que son los medios de los bordes \( AB \) y \( BF \) respectamente. ¿Cuál es la posición de las rectas \( DL \) y \( GM \)?}
{rectas secantes}{rectas paralelas}{rectas cruzantes}{}{}{}}
\End

\Data\ID{1103059403}
\Updated{2021-01-15 10:01:48}
\Level{B}
\SubArea{Lines and planes: intersecting, perpendicular, parallel}
\IMG{1103059403.pdf}
\LANGen{
\Question{
Let \( ABCDEFGH \) be a cube and \( N \) the midpoint of the edge \( EF \). See the picture. What is the mutual position of lines \( BN \) and \( CF \)?}
{skew lines}{distinct parallel lines}{intersecting lines}{}{}{}}
\LANGpl{
\Question{
Dany jest sześcian  \( ABCDEFGH \),  \( N \) to środek krawędzi  \( EF \). Zobacz rysunek. Jakie jest wzajemne położenie prostych \( BN \)  i \( CF \)?}
{proste skośne}{proste równolegle, niepokrywające się}{proste przecinające się}{}{}{}}
\LANGcs{
\Question{
Je dána krychle \( ABCDEFGH \) a bod \( N \), který je středem hrany \( EF \). Jaká je vzájemná poloha přímek \( BN \) a \( CF \)?}
{mimoběžné přímky}{různé rovnoběžné přímky}{různoběžné přímky}{}{}{}}
\LANGsk{
\Question{
Je daná kocka \( ABCDEFGH \) a bod \( N \), ktorý je stredom hrany \( EF \). Aká je vzájomná poloha priamok \( BN \) a \( CF \)?}
{mimobežné priamky}{rôzne rovnobežné priamky}{rôznobežné priamky}{}{}{}}
\LANGes{
\Question{
Sea un cubo \( ABCDEFGH \) y un punto \( N \), que es el medio del borde \( EF \). ¿Cuál es la posición de las rectas \( BN \) y \( CF \)?}
{rectas cruzantes}{rectas paralelas}{rectas secantes}{}{}{}}
\End

\Data\ID{1103059404}
\Updated{2021-01-15 10:01:07}
\Level{B}
\SubArea{Lines and planes: intersecting, perpendicular, parallel}
\IMG{1103059404.pdf}
\LANGen{
\Question{
Let \( ABCDEFGH \) be a cube and \( X \), \( Z \) the midpoints of the edges \( FB \) and \( FG \). See the picture. What is the mutual position of lines \( XZ \) and \( AH \)?}
{distinct parallel lines}{skew lines}{intersecting lines}{}{}{}}
\LANGpl{
\Question{
Dany jest sześcian \( ABCDEFGH \),  \( X \), \( Z \) to środki krawędzi  \( FB \) i \( FG \). Zobacz rysunek. Jakie jest wzajemne położenie prostych \( XZ \) i \( AH \)?}
{proste równoległe, niepokrywające się}{proste skośne}{proste przecinające się}{}{}{}}
\LANGcs{
\Question{
Je dána krychle \( ABCDEFGH \) a body \( X \), \( Z \), které jsou po řadě středy hran \( FB \) a \( FG \). Jaká je vzájemná poloha přímek \( XZ \) a \( AH \)?}
{různé rovnoběžné přímky}{mimoběžné přímky}{různoběžné přímky}{}{}{}}
\LANGsk{
\Question{
Je daná kocka \( ABCDEFGH \) a body \( X \), \( Z \), ktoré sú po rade stredy hrán \( FB \) a \( FG \). Aká je vzájomná poloha priamok \( XZ \) a \( AH \)?}
{rôzne rovnobežné priamky}{mimobežné priamky}{rôznobežné priamky}{}{}{}}
\LANGes{
\Question{
Sea un cubo \( ABCDEFGH \) y los puntos \( X \), \( Z \), que son los medios de los bordes \( FB \) y \( FG \). ¿Cuál es la posición de las rectas \( XZ \) y \( AH \)?}
{rectas paralelas}{rectas cruzantes}{rectas secantes}{}{}{}}
\End

\Data\ID{1103059405}
\Updated{2021-01-15 10:01:24}
\Level{B}
\SubArea{Lines and planes: intersecting, perpendicular, parallel}
\IMG{1103059405.pdf}
\LANGen{
\Question{
Let \( ABCDEFGH \) be a cube and \( K \), \( L \) the centers of its faces \( ABCD \) and \( BCGF \) consecutively as seen in the picture. What is the mutual position of a line \( KL \) and a plane \( CDH \)?}
{They share no common point.}{The line intersects the plane.}{The line lies on the plane.}{}{}{}}
\LANGpl{
\Question{
Dany jest sześcian \( ABCDEFGH \),  \( K \), \( L \) to środki  ścian \( ABCD \) i \( BCGF \),  zobacz rysunek. Jakie jest wzajemne położenie prostej  \( KL \) i płaszczyzny \( CDH \)?}
{Nie mają punktów wspólnych.}{Proste przecinające się.}{Prosta leży na płaszczyźnie.}{}{}{}}
\LANGcs{
\Question{
Je dána krychle \( ABCDEFGH \) a body \( K \), \( L \), které jsou po řadě středy stěn \( ABCD \) a \( BCGF \). Jaká je vzájemná poloha přímky \( KL \) a roviny \( CDH \)?}
{Přímka nemá s rovinou žádný společný bod.}{Přímka je s rovinou různoběžná.}{Přímka leží v rovině.}{}{}{}}
\LANGsk{
\Question{
Je daná kocka \( ABCDEFGH \) a body \( K \), \( L \), ktoré sú po rade stredy stien \( ABCD \) a \( BCGF \). Aká je vzájomná poloha priamky \( KL \) a roviny \( CDH \)?}
{Priamka nemá s rovinou žiadny spoločný bod.}{Priamka je s rovinou rôznobežná.}{Priamka leží v rovine.}{}{}{}}
\LANGes{
\Question{
Sea un cubo \( ABCDEFGH \) y los puntos \( K \), \( L \), que son los medios de los lados \( ABCD \) y \( BCGF \) respectamente. ¿Cuál es la posición de la recta \( KL \) y el plano \( CDH \)?}
{La recta y el plano no tienen ningún punto en común.}{La recta es secante al plano.}{La recta pertenece al plano.}{}{}{}}
\End

\Data\ID{1103059406}
\Updated{2021-01-14 05:01:56}
\Level{B}
\SubArea{Lines and planes: intersecting, perpendicular, parallel}
\IMG{1103059406.pdf}
\LANGen{
\Question{
Let \( ABCDEFGH \) be a cube and \( L \), \( N \) the centers of its faces \( BCGF \) and \( ADHE \) consecutively as seen in the picture. What is the mutual position of  the line \( LN \) and the plane \( ABG \)?}
{The line lies on the plane.}{The line intersects the plane.}{They share no common point.}{}{}{}}
\LANGpl{
\Question{
Dany jest sześcian \( ABCDEFGH \),  \( L \), \( N \) to środki  ścian \( BCGF \) i \( ADHE \), zobacz rysunek.  Jakie jest wzajemne położenie prostej \( LN \) i płaszczyzny  \( ABG \)?}
{Prosta leży na płaszczyźnie.}{Prosta przecina płaszczyznę.}{Brak punktów wspólnych.}{}{}{}}
\LANGcs{
\Question{
Je dána krychle \( ABCDEFGH \) a body \( L \), \( N \), které jsou po řadě středy stěn \( BCGF \) a \( ADHE \). Jaká je vzájemná poloha přímky \( LN \) a roviny \( ABG \)?}
{Přímka leží v rovině.}{Přímka je s rovinou různoběžná.}{Přímka nemá s rovinou žádný společný bod.}{}{}{}}
\LANGsk{
\Question{
Je daná kocka \( ABCDEFGH \) a body \( L \), \( N \), ktoré sú po rade stredy stien \( BCGF \) a \( ADHE \). Aká je vzájomná poloha priamky \( LN \) a roviny \( ABG \)?}
{Priamka leží v rovine.}{Priamka je s rovinou rôznobežná.}{Priamka nemá s rovinou žiadny spoločný bod.}{}{}{}}
\LANGes{
\Question{
Sea un cubo \( ABCDEFGH \) y los puntos \( L \), \( N \), que son los medios de los lados \( BCGF \) y \( ADHE \) respectivamente. ¿Cuál es la posición de la recta \( LN \) y el plano \( ABG \)?}
{La recta pertenece al plano.}{La recta el paralela al plano.}{La recta no tiene ningún punto en común con el plano.}{}{}{}}
\End

\Data\ID{1103059407}
\Updated{2021-01-14 05:01:09}
\Level{B}
\SubArea{Lines and planes: intersecting, perpendicular, parallel}
\IMG{1103059407.pdf}
\LANGen{
\Question{
Let \( ABCDV \) be a rectangle based-pyramid, \( V \) is its apex and \( K \), \( L \), and \( M \) are the midpoints of its edges \( AB \), \( BC \), and \( AD \) as seen in the picture. What is the mutual position of the  line \( KL \) and the plane \( CV\!M \)?}
{The line intersects the plane.}{The line lies on the plane.}{They share no common point.}{}{}{}}
\LANGpl{
\Question{
Dany jest ostrosłup\( ABCDV \), którego podstawą jest prostokąt, \( V \) jest jego wierzchołkiem a \( K \), \( L \), i \( M \) to środki krawędzi  \( AB \), \( BC \), i \( AD \); zobacz rysunek. Jakie jest wzajemne położenie prostej \( KL \) i płaszczyzny \( CV\!M \)?}
{Prosta przecina płaszczyznę.}{Prosta leży na płaszczyżnie.}{Prosta i płaszczyzna nie mają punktów wspólnych.}{}{}{}}
\LANGcs{
\Question{
Je dán pravidelný čtyřboký jehlan \( ABCDV \), kde \( V \) je hlavní vrchol jehlanu. Body \( K \), \( L \), \( M \) jsou po řadě středy hran \( AB \), \( BC \) a \( AD \). Jaká je vzájemná poloha přímky \( KL \) a roviny \( CV\!M \)?}
{Přímka je s rovinou různoběžná.}{Přímka leží v rovině.}{Přímka nemá s rovinou žádný společný bod.}{}{}{}}
\LANGsk{
\Question{
Je daný pravidelný štvorboký ihlan \( ABCDV \), kde \( V \) je hlavný vrchol ihlana. Body \( K \), \( L \), \( M \) sú po rade stredy hrán \( AB \), \( BC \) a \( AD \). Aká je vzájomná poloha priamky \( KL \) a roviny \( CV\!M \)?}
{Priamka je s rovinou rôznobežná.}{Priamka leží v rovine.}{Priamka nemá s rovinou žiadny spoločný bod.}{}{}{}}
\LANGes{
\Question{
Sea un pirámide regular cuatrilátero \( ABCDV \), donde \( V \) es el vértice del pirámide. Los puntos \( K \), \( L \), \( M \) son medios de los segmentos \( AB \), \( BC \) y \( AD \) respectivamente. ¿Cuál es la posición de la recta \( KL \) y el plano \( CV\!M \)?}
{La recta y el plano son cruzantes.}{La recta pertenece al plano.}{La recta no tiene ningún punto en común con el plano.}{}{}{}}
\End

\Data\ID{1103059408}
\Updated{2021-01-14 05:01:15}
\Level{B}
\SubArea{Lines and planes: intersecting, perpendicular, parallel}
\IMG{1103059408.pdf}
\LANGen{
\Question{
Let \( ABCDEFGH \) be a cube. See the picture. What is the mutual position of the line \( BF \) and the plane \( EGC \)?}
{They share no common point.}{The line lies on the plane.}{The line intersects the plane.}{}{}{}}
\LANGpl{
\Question{
Dany jest sześcian  \( ABCDEFGH \). Zobacz rysunek. Jakie jest wzajemne położenie prostej\( BF \) i płaszczyzny \( EGC \)?}
{Nie mają punktów wspólnych.}{Prosta leży na płaszczyźnie.}{Prosta przecina płaszczyznę.}{}{}{}}
\LANGcs{
\Question{
Je dána krychle \( ABCDEFGH \). Jaká je vzájemná poloha přímky \( BF \) a roviny \( EGC \)?}
{Přímka nemá s rovinou žádný společný bod.}{Přímka leží v rovině.}{Přímka je s rovinou různoběžná.}{}{}{}}
\LANGsk{
\Question{
Je daná kocka \( ABCDEFGH \). Aká je vzájomná poloha priamky \( BF \) a roviny \( EGC \)?}
{Priamka nemá s rovinou žiadny spoločný bod.}{Priamka leží v rovine.}{Priamka je s rovinou rôznobežná.}{}{}{}}
\LANGes{
\Question{
Sea un cubo \( ABCDEFGH \). ¿Cuál es la posición de la recta \( BF \) y el plano \( EGC \)?}
{La recta no tiene ningún punto en común con el plano.}{La recta pertenece al plano.}{La recta es secante con el plano.}{}{}{}}
\End

\Data\ID{1103059502}
\Updated{2021-01-15 10:01:52}
\Level{B}
\SubArea{Lines and planes: intersecting, perpendicular, parallel}
\IMG{1103059501.pdf}
\LANGen{
\Question{
Let \( ABCDV \) be a rectangle based-pyramid, where \( V \) is its apex and \( K \), \( L \), and \( M \) are the midpoints of its edges \( AD \), \( BC \), and \( CV \) respectively. What is the mutual position of planes \( BVK \) and \( DLM \)?}
{distinct parallel planes}{identical planes}{intersecting planes}{}{}{}}
\LANGpl{
\Question{
Dany jest ostrosłup \( ABCDV \), którego podstawą jest kwadrat, gdzie \( V \) to jego wierzchołek, natomiast  \( K \), \( L \), i \( M \) to punkty środkowe krawędzi \( AD \), \( BC \), i \( CV \).  Jakie jest wzajemne położenie płaszczyzn \( BVK \) i \( DLM \)?}
{różne płaszczyzny równoległe}{płaszczyzny identyczne}{płaszczyzny przecinające się}{}{}{}}
\LANGcs{
\Question{
Je dán pravidelný čtyřboký jehlan \( ABCDV \), kde  \( V \) je hlavní vrchol jehlanu. Body  \( K \), \( L \), \( M \) jsou po řadě středy hran \( AD \), \( BC \) a  \( CV \). Jaká je vzájemná poloha rovin \( BVK \) a \( DLM \)?}
{různé rovnoběžné roviny}{totožné roviny}{různoběžné roviny}{}{}{}}
\LANGsk{
\Question{
Je daný pravidelný štvorboký ihlan \( ABCDV \), kde  \( V \) je hlavný vrchol ihlanu. Body  \( K \), \( L \), \( M \) sú po rade stredy hrán \( AD \), \( BC \) a  \( CV \). Aká je vzájomná poloha rovín \( BVK \) a \( DLM \)?}
{rôzne rovnobežné roviny}{totožné roviny}{rôznobežné roviny}{}{}{}}
\LANGes{
\Question{
Sea un pirámide regular cuatrilátero \( ABCDV \), donde \( V \) es el vértice del pirámide. Los puntos  \( K \), \( L \), \( M \) son los medios de los bordes \( AD \), \( BC \) y \( CV \) respectivamente. ¿Cuál es la posición del los planos \( BVK \) y \( DLM \)?}
{planos paralelos}{planos coincidentes}{planos secantes}{}{}{}}
\End

\Data\ID{1103059503}
\Updated{2021-01-14 06:01:12}
\Level{B}
\SubArea{Lines and planes: intersecting, perpendicular, parallel}
\IMG{1103059503_0.pdf}
\LANGen{
\Question{
Let ABCDEFGH be a cube. What is the mutual position of planes \( ECG \), \( BDF \), and \( ABH \)?}
{three mutually intersecting planes which share one common point}{three mutually intersecting planes which share one common line}{two planes are parallel with the third intersecting them in distinct parallel lines}{}{}{}}
\LANGpl{
\Question{
Dany jest sześcian ABCDEFGH. Jakie jest wzajemne położenie płaszczyzn \( ECG \), \( BDF \), i \( ABH \)?}
{trzy wzajemnie przecinające się płaszczyzny posiadające jeden wspólny  punkt}{trzy wzajemnie przecinające się płaszczyzny posiadające jedną wspólną  prostą}{dwie płaszczyzny równoległe, trzecia płaszczyzna przecina je dwoma różnymi prostymi równoległymi}{}{}{}}
\LANGcs{
\Question{
Je dána krychle  ABCDEFGH. Jaká je vzájemná poloha tří rovin \( ECG \), \( BDF \) a \( ABH \)?}
{tři vzájemně různoběžné roviny se společným jediným bodem}{tři vzájemně různoběžné roviny se společnou jedinou přímkou}{dvě roviny jsou rovnoběžné a třetí je protíná v různých rovnoběžných přímkách}{}{}{}}
\LANGsk{
\Question{
Je daná kocka  ABCDEFGH. Aká je vzájomná poloha troch rovín \( ECG \), \( BDF \) a \( ABH \)?}
{tri vzájomne rôznobežné roviny so spoločným jediným bodom}{tri vzájomne rôznobežné roviny so spoločnou jedinou priamkou}{dve roviny sú rovnobežné a tretia ich pretína v rôznych rovnobežných priamkách}{}{}{}}
\LANGes{
\Question{
Sea un cubo \( ABCDEFGH \). ¿Cuál es la posición de los tres planos \( ECG \), \( BDF \) y \( ABH \)?}
{tres planos secantes con un único punto en común}{tres planos secantes son una recta en común}{dos planos son paralelos y el tercero los seca}{}{}{}}
\End

\Data\ID{1103059504}
\Updated{2021-01-15 10:01:19}
\Level{B}
\SubArea{Lines and planes: intersecting, perpendicular, parallel}
\IMG{1103059503.pdf}
\LANGen{
\Question{
Let \( ABCDEFGH \) be a cube with \( K \), and \( L \) being the midpoints of its edges \( AE \) and \( CG \), and let \( M \) be the centre of its face \( ABFE \). What is the mutual position of planes \( BCE \), \( ADF \), and \( KLM \)?}
{three mutually intersecting planes which share one common line}{three mutually intersecting planes which share one common point}{two planes are parallel with the third intersecting them in distinct parallel lines}{}{}{}}
\LANGpl{
\Question{
Dany jest sześcian \( ABCDEFGH \), gdzie \( K \), i  \( L \) to punkty środkowe krawędzi  \( AE \) i \( CG \), punkt \( M \) to punkt środkowy ściany  \( ABFE \). Jakie jest wzajemne położenie płaszczyzn  \( BCE \), \( ADF \), i \( KLM \)?}
{trzy wzajemnie przecinające się płaszczyzny posiadające jedną wspólną  prostą}{trzy wzajemnie przecinające się płaszczyzny posiadające jeden wspólny  punkt}{dwie równoległe płaszczyzny oraz trzecia przecinająca je niepokrywającymi prostymi równoległymi}{}{}{}}
\LANGcs{
\Question{
Je dána krychle \( ABCDEFGH \). Body  \( K \), \( L \)  jsou po řadě středy hran \( AE \), \( CG \) a bod \( M \) je středem stěny \( ABFE \). Jaká je vzájemná poloha tří rovin \( BCE \), \( ADF \) a \( KLM \)?}
{tři vzájemně různoběžné roviny se společnou jedinou přímkou}{tři vzájemně různoběžné roviny se společným jediným bodem}{dvě roviny jsou rovnoběžné a třetí je protíná v různých rovnoběžných přímkách}{}{}{}}
\LANGsk{
\Question{
Je daná kocka \( ABCDEFGH \). Body  \( K \), \( L \)  sú po rade stredy hrán \( AE \), \( CG \) a bod \( M \) je stredom steny \( ABFE \). Aká je vzájomná poloha troch rovín \( BCE \), \( ADF \) a \( KLM \)?}
{tri vzájomne rôznobežné roviny so spoločnou jedinou priamkou}{tri vzájomne rôznobežné roviny so spoločným jediným bodom}{dve roviny sú rovnobežné a tretia ich pretína v rôznych rovnobežných priamkach}{}{}{}}
\LANGes{
\Question{
Sea un cubo \( ABCDEFGH \). Los puntos \( K \), \( L \)  son los medios de los bordes \( AE \) y \( CG \) respectivamente y el punto \( M \) es el medio del lado \( ABFE \). ¿Cuál es la posición de los tres planos \( BCE \), \( ADF \) y \( KLM \)?}
{tres planos secantes con una recta en común}{tres planos secantes con un único punto en común}{dos planos son paralelos y el tercero los seca en rectas paralelas distintas}{}{}{}}
\End

\Data\ID{1103059505}
\Updated{2021-01-15 10:01:46}
\Level{B}
\SubArea{Lines and planes: intersecting, perpendicular, parallel}
\IMG{1103059504.pdf}
\LANGen{
\Question{
Let \( ABCDEFGH \) be a cube and \( X \) be the midpoint of its edge \( AE \). What is the cross-section of the cube if we slice it with a plane \( BGX \)?}
{a quadrilateral \( BGPX \) with \( P \) being the midpoint of the edge \( EH \)}{a quadrilateral \( BGHX \)}{a triangle \( BGX \)}{a quadrilateral \( BGPX \) with \( P \) being the midpoint of the edge \( DH \)}{}{}}
\LANGpl{
\Question{
Dany jest sześcian \( ABCDEFGH \), gdzie  \( X \) to punkt środkowy krawędzi  \( AE \). Jaki jest przekrój  sześcianu jeśli  dokonano go płaszczyzną \( BGX \)?}
{czworokąt  \( BGPX \), gdzie \( P \) to punkt środkowy krawędzi  \( EH \)}{czworokąt \( BGHX \)}{trójkąt  \( BGX \)}{czworokąt \( BGPX \), gdzie \( P \) to punkt środkowy krawędzi  \( DH \)}{}{}}
\LANGcs{
\Question{
Je dána krychle \( ABCDEFGH \). Bod \( X \) je středem hrany \( AE \). Řezem dané krychle rovinou \( BGX \) je:}
{čtyřúhelník \( BGPX \), kde bod \( P \) je střed hrany \( EH \)}{čtyřúhelník \( BGHX \)}{trojúhelník \( BGX \)}{čtyřúhelník \( BGPX \), kde bod \( P \) je střed hrany \( DH \)}{}{}}
\LANGsk{
\Question{
Je daná kocka \( ABCDEFGH \). Bod \( X \) je stredom hrany \( AE \). Rezom danej kocky rovinou \( BGX \) je:}
{štvoruholník \( BGPX \), kde bod \( P \) je stred hrany \( EH \)}{štvoruholník \( BGHX \)}{trojuholník \( BGX \)}{štvoruholník \( BGPX \), kde bod \( P \) je stred hrany \( DH \)}{}{}}
\LANGes{
\Question{
Sea un cubo \( ABCDEFGH \). El punto \( X \) es el medio del borde \( AE \). El corte del cubo por el plano \( BGX \) es:}
{cuatrilátero \( BGPX \), donde el punto \( P \) es el medio del borde \( EH \)}{cuatrilátero \( BGHX \)}{triángulo \( BGX \)}{cuatrilátero \( BGPX \), donde el punto \( P \) es el medio del borde \( DH \)}{}{}}
\End

\Data\ID{1103059506}
\Updated{2021-01-15 10:01:20}
\Level{B}
\SubArea{Lines and planes: intersecting, perpendicular, parallel}
\IMG{1103059505.pdf}
\LANGen{
\Question{
Let \( ABCDEFGH \) be a cube and let \( X \), \( Y \), and \( Z \) be the midpoints of edges \( AB \), \( AE \), and \( CG \) respectively. What is the cross-section of the cube if we slice it with a plane \( XYZ \)?}
{a hexagon \( XLZKMY \) with points \( L \), \( K \), and \( M \) lying on edges \( BC \), \( GH \), and \( EH \) respectively}{a pentagon \( XLZKY \) with points \( L \) and \( K \) lying on edges \( BC \) and \( GH \) respectively}{a triangle \( XYZ \)}{a quadrilateral \( XZKY \) with \( K \) being the midpoint of the edge \( GH \)}{}{}}
\LANGpl{
\Question{
Dany jest sześcian \( ABCDEFGH \), gdzie  \( X \), \( Y \), i \( Z \) to punkty środkowe krawędzi  \( AB \), \( AE \), i \( CG \). Jaki jest przekrój sześcianu jeśli dokonano go płaszczyzną  \( XYZ \)?}
{sześciokąt  \( XLZKMY \), którego punkty  \( L \), \( K \), i \( M \) leżą odpowiednio na krawędziach  \( BC \), \( GH \), i \( EH \)}{pięciokąt  \( XLZKY \), którego punkty \( L \) i \( K \) leżą odpowiednio na krawędziach \( BC \) i \( GH \)}{trójkąt \( XYZ \)}{czworokąt  \( XZKY \), gdzie  \( K \) to punkt środkowy krawędzi  \( GH \)}{}{}}
\LANGcs{
\Question{
Je dána krychle \( ABCDEFGH \). Body \( X \), \( Y \), \( Z \) jsou po řadě středy hran \( AB \), \( AE \) a \( CG \). Řezem dané krychle rovinou \( XYZ \) je:}
{šestiúhelník \( XLZKMY \), kde body  \( L \), \( K \), \( M \) leží po řadě na hranách \( BC \), \( GH \) a \( EH \)}{pětiúhelník \( XLZKY \), kde body \( L \), \( K \) leží po řadě na hranách \( BC \) a \( GH \)}{trojúhelník \( XYZ \)}{čtyřúhelník \( XZKY \), kde bod \( K \) je střed hrany \( GH \)}{}{}}
\LANGsk{
\Question{
Je daná kocka \( ABCDEFGH \). Body \( X \), \( Y \), \( Z \) sú po rade stredy hrán \( AB \), \( AE \) a \( CG \). Rezom danej kocky rovinou \( XYZ \) je:}
{šesťuholník \( XLZKMY \), kde body  \( L \), \( K \), \( M \) ležia po rade na hranách \( BC \), \( GH \) a \( EH \)}{päťuholník \( XLZKY \), kde body \( L \), \( K \) ležia po rade na hranách \( BC \) a \( GH \)}{trojuholník \( XYZ \)}{štvoruholník \( XZKY \), kde bod \( K \) je stred hrany \( GH \)}{}{}}
\LANGes{
\Question{
Sea un cubo \( ABCDEFGH \). Los puntos \( X \), \( Y \), \( Z \) son los medios de los bordes \( AB \), \( AE \) y \( CG \) respectivamente. El corte del cubo por el plano \( XYZ \) es:}
{hexágono \( XLZKMY \), donde los puntos  \( L \), \( K \), \( M \) pertenecen a los bordes \( BC \), \( GH \) y \( EH \) respectivamente}{pentágono\( XLZKY \), donde los puntos \( L \), \( K \) pertenecen a los bordes \( BC \) a \( GH \) respectivamente}{triángulo \( XYZ \)}{cuadrilátero \( XZKY \), donde el punto \( K \) es el medio del borde \( GH \)}{}{}}
\End

\Data\ID{1103059507}
\Updated{2021-01-15 10:01:52}
\Level{B}
\SubArea{Lines and planes: intersecting, perpendicular, parallel}
\IMG{1103059506.pdf}
\LANGen{
\Question{
Let \( ABCDEFGH \) be a cube with \( K \) and \( L \) being the midpoints of edges \( AE \) and \( AB \) respectively, and let \( M \) be the midpoint of the face diagonal \( EG \). What is the cross-section of the cube if we slice it with a plane \( KLM \)?}
{a pentagon \( KLPQR \) with points \( P \), \( Q \), and \( R \) lying on edges \( BC \), \( FG \), and \( EH \) respectively}{a triangle \( KLM \)}{a pentagon \( KLPQM \) with points \( P \) and \( Q \) lying on edges \( BC \) and \( FG \) respectively}{a quadrilateral \( KLMR \) with point \( R \) lying on the  edge \( EH \)}{}{}}
\LANGpl{
\Question{
Dany jest sześcian \( ABCDEFGH \), gdzie  \( K \) i \( L \) to odpowiednio punkty środkowe krawędzi  \( AE \) i \( AB \), punkt  \( M \) to punkt środkowy przekątnej  \( EG \). Jaki jest przekrój sześcianu, którego dokonano płaszczyzną \( KLM \)?}
{pięciokąt  \( KLPQR \), gdzie punkty  \( P \), \( Q \), i \( R \) leżą odpowiednio na krawędziach \( BC \), \( FG \), i \( EH \)}{trójkąt  \( KLM \)}{pięciokąt  \( KLPQM \), gdzie punkty  \( P \) i \( Q \) leżą odpowiednio na krawędziach \( BC \) i \( FG \)}{czworokąt  \( KLMR \), gdzie punkt  \( R \) leży  na krawędzi  \( EH \)}{}{}}
\LANGcs{
\Question{
Je dána krychle \( ABCDEFGH \). Body \( K \), \( L \) jsou po řadě středy hran \( AE \) a \( AB \). Bod \( M \) je střed stěnové úhlopříčky \( EG \). Řezem dané krychle rovinou \( KLM \) je:}
{pětiúhelník \( KLPQR \), kde body \( P \), \( Q \), \( R \)  leží po řadě na hranách \( BC \), \( FG \) a \( EH \)}{trojúhelník \( KLM \)}{pětiúhelník \( KLPQM \), kde body \( P \), \( Q \) leží po řadě na hranách \( BC \) a \( FG \)}{čtyřúhelník \( KLMR \), kde bod \( R \) leží na hraně \( EH \)}{}{}}
\LANGsk{
\Question{
Je daná kocka \( ABCDEFGH \). Body \( K \), \( L \) sú po rade stredy hrán \( AE \) a \( AB \). Bod \( M \) je stred stenovej uhlopriečky \( EG \). Rezom danej kocky rovinou \( KLM \) je:}
{päťuholník \( KLPQR \), kde body \( P \), \( Q \), \( R \)  leží po rade na hranách \( BC \), \( FG \) a \( EH \)}{trojuholník \( KLM \)}{päťuholník \( KLPQM \), kde body \( P \), \( Q \) leží po rade na hranách \( BC \) a \( FG \)}{štvoruholník \( KLMR \), kde bod \( R \) leží na hrane \( EH \)}{}{}}
\LANGes{
\Question{
Sea un cubo \( ABCDEFGH \). Los puntos \( K \), \( L \) son los medios de los bordes \( AE \) y \( AB \) respectivamente. El punto \( M \) es el medio de la diagonal \( EG \). El corte del cubo por el plano \( KLM \) es:}
{pentágono \( KLPQR \), donde los puntos \( P \), \( Q \), \( R \)  pertenecen a los bordes \( BC \), \( FG \) y \( EH \) respectivamente}{triángulo \( KLM \)}{pentágono \( KLPQM \), donde los puntos \( P \), \( Q \) pertenecen a los bordes \( BC \) y \( FG \) respectivamente}{cuadrilátero \( KLMR \), donde el punto \( R \) pertenece al borde \( EH \)}{}{}}
\End

\Data\ID{1103059601}
\Updated{2021-01-15 09:01:42}
\Level{B}
\SubArea{Lines and planes: intersecting, perpendicular, parallel}
\IMG{1103059601.pdf}
\LANGen{
\Question{
Let \( ABCDV \) be a rectangle based-pyramid and \( V \) its apex. The pyramid is sliced by a plane \( EFG \) which is defined by:
\begin{align*}
 E&\in BC\ \wedge\ |BE|=2|CE|, \\
F&\in AV\ \wedge\ |AF|=2|VF|, \\
G&\in DV\ \wedge\ |DG|=2|VG|
\end{align*}
(see the picture). What is the cross-section of the pyramid  if we slice it with the plane \( EFG \)?}
{a trapezium \( BCGF \)}{a triangle \( EFG \)}{a triangle \( AEV \)}{a pentagon \( ABEGF \)}{}{}}
\LANGpl{
\Question{
Dany jest ostrosłup \( ABCDV \), którego podstawą jest prostokąt, punkt \( V \) to jego wierzchołek. Ostrosłup przekrojono płaszczyzną \( EFG \) określoną przez:
\begin{align*}
 E&\in BC\ \wedge\ |BE|=2|CE|, \\
F&\in AV\ \wedge\ |AF|=2|VF|, \\
G&\in DV\ \wedge\ |DG|=2|VG|
\end{align*}
(spójrz na rysunek). Jaki jest przekrój ostrosłupa, którego dokonano płaszczyzną  \( EFG \)?}
{trapez \( BCGF \)}{trójkąt \( EFG \)}{trójkąt \( AEV \)}{pięciokąt \( ABEGF \)}{}{}}
\LANGcs{
\Question{
Je dán pravidelný čtyřboký jehlan \( ABCDV \), kde  \( V \) je hlavní vrchol jehlanu. Rovina řezu \( EFG \) je určena takto:
\begin{align*}
 E&\in BC\ \wedge\ |BE|=2|CE|, \\
F&\in AV\ \wedge\ |AF|=2|VF|, \\
G&\in DV\ \wedge\ |DG|=2|VG|
\end{align*}]
(viz obrázek). Řezem daného jehlanu rovinou \( EFG \) je:}
{lichoběžník \( BCGF \)}{trojúhelník \( EFG \)}{trojúhelník \( AEV \)}{pětiúhelník \( ABEGF \)}{}{}}
\LANGsk{
\Question{
Je daný pravidelný štvorboký ihlan \( ABCDV \), kde  \( V \) je hlavný vrchol ihlanu. Rovina rezu \( EFG \) je určená takto:
\begin{align*}
 E&\in BC\ \wedge\ |BE|=2|CE|, \\
F&\in AV\ \wedge\ |AF|=2|VF|, \\
G&\in DV\ \wedge\ |DG|=2|VG|
\end{align*}
(viď obrázok). Rezom daného ihlanu rovinou \( EFG \) je:}
{lichobežník \( BCGF \)}{trojuholník \( EFG \)}{trojuholník \( AEV \)}{päťuholník \( ABEGF \).}{}{}}
\LANGes{
\Question{
Sea un pirámide regular cuadrilátero \( ABCDV \), donde  \( V \) es el vértice del pirámide. El plano del corte \( EFG \) es definido:
\begin{align*}
 E&\in BC\ \wedge\ |BE|=2|CE|, \\
F&\in AV\ \wedge\ |AF|=2|VF|, \\
G&\in DV\ \wedge\ |DG|=2|VG|
\end{align*}]
(vea el dibujo). El corte del pirámide por el plano \( EFG \) es:}
{trapezoide \( BCGF \)}{triángulo \( EFG \)}{triángulo \( AEV \)}{pentágono \( ABEGF \)}{}{}}
\End

\Data\ID{1103059602}
\Updated{2021-01-15 09:01:37}
\Level{B}
\SubArea{Lines and planes: intersecting, perpendicular, parallel}
\IMG{1103059602.pdf}
\LANGen{
\Question{
Let \( ABCDV \) be a rectangle based-pyramid and \( V \) its apex. The pyramid is sliced by a plane \( XYZ \) which is defined by: 
\begin{align*}
X&\text{ is the midpoint of  the edge }AD,\\
Y&\in CD\ \wedge\ |DY|=3|CY|,\\
Z&\in BV\ \wedge\ |BZ|=3|VZ|
\end{align*}
(see the picture). What is the cross-section of the pyramid if we slice it with the plane \( XYZ \)?}
{a pentagon \( XYKZL \) with points \( K \) and \( L \) lying on the edges \( CV \) and \( AV \)}{a triangle \( XYZ \)}{a quadrilateral \( XYZL \) with point \( L \) lying on the edge \( AV \)}{a quadrilateral \( XYKZ \) with point \( K \) lying on the edge \( CV \)}{}{}}
\LANGpl{
\Question{
Dany jest ostrosłup \( ABCDV \), którego podstawą jest prostokąt, punkt \( V \) to jego wierzchołek. Ostrosłup przekrojono płaszczyzną  \( XYZ \) określoną przez: 
\begin{align*}
X&\text{ to punkt środkowy krawędzi }AD, \\
Y&\in CD\ \wedge\ |DY|=3|CY|, \\ 
Z&\in BV\ \wedge\ |BZ|=3|VZ| 
\end{align*}
(spójrz na rysunek). Jaki jest przekrój ostrosłupa jeśli dokonano go płaszczyzną \( XYZ \)?}
{pięciokąt  \( XYKZL \), gdzie punkty \( K \) i \( L \) leżą na krawędziach  \( CV \) i \( AV \)}{trójkąt \( XYZ \)}{czworokąt  \( XYZL \), gdzie punkt \( L \) leży na krawędzi  \( AV \)}{czworokąt  \( XYKZ \), gdzie punkt \( K \) leży na krawędzi  \( CV \)}{}{}}
\LANGcs{
\Question{
Je dán pravidelný čtyřboký jehlan \( ABCDV \), kde  \( V \) je hlavní vrchol jehlanu. Rovina řezu \( XYZ \) je určena takto: 
\begin{align*}
X&\text{ je střed hrany }AD,\\
Y&\in CD\ \wedge\ |DY|=3|CY|, \\
Z&\in BV\ \wedge\ |BZ|=3|VZ|
\end{align*}
(viz obrázek). Řezem daného jehlanu rovinou \( XYZ \) je:}
{pětiúhelník \( XYKZL \), kde body \( K \) a \( L \) leží po řadě na hranách  \( CV \) a \( AV \)}{trojúhelník \( XYZ \)}{čtyřúhelník \( XYZL \), kde bod \( L \) leží na hraně \( AV \)}{čtyřúhelník \( XYKZ \), kde bod \( K \) leží na hraně \( CV \)}{}{}}
\LANGsk{
\Question{
Je daný pravidelný štvorboký ihlan \( ABCDV \), kde  \( V \) je hlavný vrchol ihlanu. Rovina rezu \( XYZ \) je určená takto: 
 \begin{align*}
X&\text{ je stred hrany }AD, \\
Y&\in CD\ \wedge\ |DY|=3|CY|,\\
Z&\in BV\ \wedge\ |BZ|=3|VZ|
\end{align*}
(viď obrázok). Rezom daného ihlanu rovinou \( XYZ \) je:}
{päťuholník \( XYKZL \), kde body \( K \) a \( L \) leží po rade na hranách  \( CV \) a \( AV \)}{trojuholník \( XYZ \)}{štvoruholník \( XYZL \), kde bod \( L \) leží na hrane \( AV \)}{štvoruholník \( XYKZ \), kde bod \( K \) leží na hrane \( CV \)}{}{}}
\LANGes{
\Question{
Sea un pirámide regular cuadrilátero \( ABCDV \), donde  \( V \) es el vértice del pirámide. El plano de corte \( XYZ \) se define:
\begin{align*}
X&\text{ es el medio del borde }AD,\\
Y&\in CD\ \wedge\ |DY|=3|CY|, \\
Z&\in BV\ \wedge\ |BZ|=3|VZ|
\end{align*}
(vea el dibujo). El corte del pirámide por el plano \( XYZ \) es:}
{el pentágono \( XYKZL \), donde los puntos \( K \) y \( L \) pertenecen a los bordes \( CV \) y \( AV \) respectivamente}{triángulo \( XYZ \)}{cuadrado \( XYZL \), donde el punto \( L \) pertenece al borde \( AV \)}{cuadrado \( XYKZ \), donde el punto \( K \) pertenece al borde \( CV \)}{}{}}
\End

\Data\ID{1103059603}
\Updated{2021-01-15 10:01:17}
\Level{B}
\SubArea{Lines and planes: intersecting, perpendicular, parallel}
\IMG{1103059603.pdf}
\LANGen{
\Question{
Let \( ABCDEFGH \) be a cube and let \( XY \) be a line where:
\begin{align*}
X&\text{ lays on a ray }BC\text{ and }|BX|=1.5|BC|,\\
Y&\text{ lays on a ray }HE\text{ and }|HY|=1.5|HE|
\end{align*}
(see the picture). The points of intersection of the line \( XY \) with the surface of the cube lay:}
{on the sides \( ABFE \) and \( DCGH \)}{on the side \( ABFE \) and the edge \( CG \)}{on the edges \( AE \) and \( CG \)}{on the sides \( ADHE \) and \( BCGF \)}{}{}}
\LANGpl{
\Question{
Dany jest sześcian \( ABCDEFGH \), prosta \( XY \) to prosta, gdzie:
\begin{align*}
X&\text{ leży na półprostej  }BC\text{ i }|BX|=1{,}5|BC|,\\
Y&\text{ leży na półprostej }HE\text{ i }|HY|=1{,}5|HE|
\end{align*}
(spójrz na rysunek). Punkty przecięcia prostej \( XY \) z powierzchnią sześcianu leżą:}
{na ścianach   \( ABFE \) i \( DCGH \)}{na ścianie \( ABFE \) oraz krawędzi  \( CG \)}{na krawędziach \( AE \) i \( CG \)}{na ścianach \( ADHE \) i \( BCGF \)}{}{}}
\LANGcs{
\Question{
Je dána krychle \( ABCDEFGH \) a přímka \( XY \), která je určena takto:
\begin{align*}
X&\text{ leží na polopřímce }BC\text{ a }|BX|=1{,}5|BC|,\\
Y&\text{ leží na polopřímce }HE\text{ a }|HY|=1{,}5|HE|
\end{align*}
(viz obrázek). Průsečíky přímky \( XY \) s povrchem krychle leží:}
{ve stěnách krychle \( ABFE \) a \( DCGH \)}{ve stěně krychle \( ABFE \) a na hraně \( CG \)}{na hranách krychle \( AE \) a \( CG \)}{ve stěnách krychle \( ADHE \) a \( BCGF \)}{}{}}
\LANGsk{
\Question{
Je daná kocka \( ABCDEFGH \) a priamka \( XY \), ktorá je určená takto:
\begin{align*}
X&\text{ leží na polpriamke }BC\text{ a }|BX|=1{,}5|BC|,\\
Y&\text{ leží na polpriamke }HE\text{ a }|HY|=1{,}5|HE|
\end{align*}
(viď obrázok). Priesečníky priamky \( XY \) s povrchom kocky leží:}
{v stenách kocky \( ABFE \) a \( DCGH \)}{v stene kocky \( ABFE \) a na hrane \( CG \)}{na hranách kocky \( AE \) a \( CG \)}{v stenách kocky \( ADHE \) a \( BCGF \)}{}{}}
\LANGes{
\Question{
Sea un cubo \( ABCDEFGH \) y la recta \( XY \), definida:
\begin{align*}
X&\text{ pertenece a la semirecta }BC\text{ a }|BX|=1{,}5|BC|,\\
Y&\text{ pertenece a la semirecta }HE\text{ a }|HY|=1{,}5|HE|
\end{align*}
(vea el dibujo). Las intersecciones de la recta \( XY \) con el cubo pertenecen a:}
{los lados del cubo \( ABFE \) y \( DCGH \)}{al lado \( ABFE \) y al borde \( CG \)}{a los bordes \( AE \) y \( CG \)}{a los lados \( ADHE \) y \( BCGF \)}{}{}}
\End

\Data\ID{1103059604}
\Updated{2021-01-15 10:01:47}
\Level{B}
\SubArea{Lines and planes: intersecting, perpendicular, parallel}
\IMG{1103059604.pdf}
\LANGen{
\Question{
Let \( ABCDEFGH \) be a cube and let \( XY \) be a line where:
\begin{align*}
X&\text{ lays on a ray }DH\text{ and }|DX|=1.5|DH|,\\
Y&\text{ lays on a ray }DB\text{ and }|DB|=|BY|
\end{align*}
(see the picture). The points of intersection of the line \( XY \) with the surface of the cube lay:}
{on the side \( EFGH \) and the edge \( BF \)}{on the edges \( EF \) and \( BF \)}{on the sides \( EFGH \) and \( ABCD \)}{on the edges \( HG \) and \( BF \)}{}{}}
\LANGpl{
\Question{
Dany jest sześcian \( ABCDEFGH \) oraz prosta \( XY \), gdzie:
\begin{align*}
X&\text{ leży na półprostej  }DH\text{ i }|DX|=1{,}5|DH|,\\
Y&\text{ leży na półprostej }DB\text{ i }|DB|=|BY|
\end{align*}
(spójrz na rysunek). Punkty przecięcia prostej \( XY \) z powierzchnią sześcianu leżą:}
{na ścianie \( EFGH \) oraz krawędzi \( BF \)}{na krawędziach  \( EF \) i \( BF \)}{na ścianach \( EFGH \) i \( ABCD \)}{na krawędziach \( HG \) i \( BF \)}{}{}}
\LANGcs{
\Question{
Je dána krychle \( ABCDEFGH \) a přímka \( XY \), která je určena takto:
\begin{align*}
X&\text{ leží na polopřímce }DH\text{ a }|DX|=1{,}5|DH|,\\
Y&\text{ leží na polopřímce }DB\text{ a }|DB|=|BY|
\end{align*}
(viz obrázek). Průsečíky přímky \( XY \) s povrchem krychle leží:}
{ve stěně krychle \( EFGH \) a na hraně \( BF \)}{na hranách krychle \( EF \) a \( BF \)}{ve stěnách krychle \( EFGH \) a \( ABCD \)}{na hranách krychle \( HG \) a \( BF \)}{}{}}
\LANGsk{
\Question{
Je daná kocka \( ABCDEFGH \) a priamka \( XY \), ktorá je určená takto:
\begin{align*}
X&\text{ leží na polpriamke }DH\text{ a }|DX|=1{,}5|DH|,\\
Y&\text{ leží na polpriamke }DB\text{ a }|DB|=|BY|
\end{align*}
(viď obrázok). Priesečníky priamky \( XY \) s povrchom kocky leží:}
{v stene kocky \( EFGH \) a na hrane \( BF \)}{na hranách kocky \( EF \) a \( BF \)}{v stenách kocky \( EFGH \) a \( ABCD \)}{na hranách kocky \( HG \) a \( BF \)}{}{}}
\LANGes{
\Question{
Sea un cubo \( ABCDEFGH \) y una recta \( XY \), definida:
\begin{align*}
X&\text{ pertenece a la semirecta }DH\text{ y }|DX|=1{,}5|DH|,\\
Y&\text{ pertenece a la semirecta }DB\text{ y }|DB|=|BY|
\end{align*}
(vea el dibujo). Intersecciones de la recta \( XY \) con la superficie del cubo pertenecen:}
{al lado \( EFGH \) y al borde \( BF \)}{a los bordes \( EF \) y \( BF \)}{a los lados \( EFGH \) y \( ABCD \)}{a los bordes \( HG \) y \( BF \)}{}{}}
\End

\Data\ID{1103059605}
\Updated{2021-01-15 10:01:32}
\Level{B}
\SubArea{Lines and planes: intersecting, perpendicular, parallel}
\IMG{1103059605_0.pdf}
\LANGen{
\Question{
Let \( ABCDEFGH \) be a cube and let \( XY \) be a line where:
\begin{align*}
X&\text{ lays on a ray }CB\text{ and }|CX|=1.5|BC|,\\
Y&\text{ lays on a ray }EH\text{ and }|EY|=1.5|EH|
\end{align*}
(see the picture). The points of intersection of the line \( XY \) with the surface of the cube lay:}
{on the sides \( ABFE \) and \( DCGH \)}{on the sides \( EFGH \) and \( ABCD \)}{on the side \( ABCD \) and the edge \( HG \)}{on the edges \( HG \) and \( AB \)}{}{}}
\LANGpl{
\Question{
Dany jest sześcian \( ABCDEFGH \) oraz prosta \( XY \), gdzie:
\begin{align*}
X&\text{ leży na półprostej  }CB\text{ i }|CX|=1{,}5|BC|,\\
Y&\text{ leży na półprostej }EH\text{ i }|EY|=1{,}5|EH|
\end{align*}
(spójrz na rysunek). Punkty przecięcia prostej  \( XY \) z powierzchnią sześcianu leżą:}
{na ścianach \( ABFE \) i \( DCGH \)}{na ścianach \( EFGH \) i \( ABCD \)}{na ścianie \( ABCD \) oraz krawędzi  \( HG \)}{na krawędziach \( HG \) i \( AB \)}{}{}}
\LANGcs{
\Question{
Je dána krychle \( ABCDEFGH \) a přímka \( XY \), která je určena takto:
\begin{align*}
X&\text{ leží na polopřímce }CB\text{ a }|CX|=1{,}5|BC|,\\
Y&\text{ leží na polopřímce }EH\text{ a }|EY|=1{,}5|EH|
\end{align*}
(viz obrázek). Průsečíky přímky \( XY \) s povrchem krychle leží:}
{ve stěnách krychle \( ABFE \) a \( DCGH \)}{ve stěnách krychle \( EFGH \) a \( ABCD \)}{ve stěně krychle \( ABCD \) a na hraně \( HG \)}{na hranách krychle \( HG \) a \( AB \)}{}{}}
\LANGsk{
\Question{
Je daná kocka \( ABCDEFGH \) a priamka \( XY \), ktorá je určená takto:
\begin{align*}
X&\text{ leží na polpriamke }CB\text{ a }|CX|=1{,}5|BC|,\\
Y&\text{ leží na polpriamke }EH\text{ a }|EY|=1{,}5|EH| 
\end{align*}
(viď obrázok). Priesečníky priamky \( XY \) s povrchom kocky leží:}
{v stenách kocky \( ABFE \) a \( DCGH \)}{v stenách kocky \( EFGH \) a \( ABCD \)}{v stene kocky \( ABCD \) a na hrane \( HG \)}{na hranách kocky \( HG \) a \( AB \)}{}{}}
\LANGes{
\Question{
Sea un cubo \( ABCDEFGH \) y una recta \( XY \), definida:
\begin{align*}
X&\text{ pertenece a la semirecta }CB\text{ y }|CX|=1{,}5|BC|,\\
Y&\text{ pertenece a la semirecta }EH\text{ y }|EY|=1{,}5|EH|
\end{align*}
(vea el dibujo). Intersecciones de la recta \( XY \) con la superficie del cubo pertenecen:}
{a los lados \( ABFE \) y \( DCGH \)}{a los lados \( EFGH \) y \( ABCD \)}{al lado \( ABCD \) y al borde \( HG \)}{a los bordes \( HG \) y \( AB \)}{}{}}
\End

\Data\ID{1103059606}
\Updated{2021-01-15 09:01:43}
\Level{B}
\SubArea{Lines and planes: intersecting, perpendicular, parallel}
\IMG{1103059606.pdf}
\LANGen{
\Question{
Let \( ABCDV \) be a rectangle based-pyramid, where \( V \) is its apex, and let \( XY \) be a line where:
\begin{align*}
X&\text{ lays on the side }AV\text{ and }|AX|=|XV|,\\
Y&\text{ lays on a ray }DC\text{ and }|DY|=1.5|DC|
\end{align*}
(see the picture). The points of intersection of the line \( XY \) with the surface of the pyramid are:}
{the point \( X \) and a point on the pyramid face \( BCV \)}{the point \( X \) and a point on the pyramid face \( DCV \)}{the point \( X \) and a point on the pyramid edge \( CV \)}{the point \( X \) only}{}{}}
\LANGpl{
\Question{
Dany jest ostrosłup \( ABCDV \), którego podstawą jest prostokąt, gdzie punkt \( V \) to jego wierzchołek, prosta \( XY \) to prosta, gdzie:
\begin{align*}
X&\text{ leży na ścianie }AV\text{ i }|AX|=|XV|,\\
Y&\text{ leży na półprostej }DC\text{ i }|DY|=1{,}5|DC|
\end{align*}
(spójrz na rysunek). Punkty przecięcia prostej \( XY \) z powierzchnią ostrosłupa to:}
{punkt \( X \) oraz punkt na ścianie \( BCV \)}{punkt \( X \) oraz punkt na ścianie \( DCV \)}{punkt \( X \) oraz punkt na krawędzi \( CV \)}{tylko punkt \( X \)}{}{}}
\LANGcs{
\Question{
Je dán pravidelný čtyřboký jehlan \( ABCDV \), kde \( V \) je hlavní vrchol jehlanu. Přímka \( XY \) je určena takto:
\begin{align*}
X&\text{ leží na hraně }AV\text{ a }|AX|=|XV|,\\
Y&\text{ leží na polopřímce }DC\text{ a }|DY|=1{,}5|DC|
\end{align*}
(viz obrázek). Průsečíky přímky \( XY \) s povrchem jehlanu jsou:}
{bod \( X \) a bod ležící ve stěně jehlanu \( BCV \)}{bod \( X \) a bod ležící ve stěně jehlanu \( DCV \)}{bod  \( X \) a bod ležící na hraně jehlanu \( CV \)}{pouze bod \( X \)}{}{}}
\LANGsk{
\Question{
Je daný pravidelný štvorboký ihlan \( ABCDV \), kde \( V \) je hlavný vrchol ihlana. Priamka \( XY \) je určená takto:
\begin{align*}
X&\text{ leží na hrane }AV\text{ a }|AX|=|XV|,\\
Y&\text{ leží na polpriamke }DC\text{ a }|DY|=1{,}5|DC|
\end{align*}
(viď obrázok). Priesečníky priamky \( XY \) s povrchom ihlanu sú:}
{bod \( X \) a bod ležiaci v stene ihlanu \( BCV \)}{bod \( X \) a bod ležiaci v stene ihlanu \( DCV \)}{bod  \( X \) a bod ležiaci na hrane ihlanu \( CV \)}{len bod \( X \)}{}{}}
\LANGes{
\Question{
Sea un pirámide regular cuadrilátero \( ABCDV \), donde  \( V \) es el vértice del pirámide. La recta \( XY \) es definida:
\begin{align*}
X&\text{ pertenece al borde }AV\text{ y }|AX|=|XV|,\\
Y&\text{ pertenece a la semirecta }DC\text{ y }|DY|=1{,}5|DC|
\end{align*}
(vea el dibujo). Intersecciones de la recta \( XY \) con la superficie del pirámide son:}
{el punto \( X \) y un punto perteneciente al lado \( BCV \)}{el punto \( X \) y un punto perteneciente al lado \( DCV \)}{el punto \( X \) y un punto perteneciente al borde \( CV \)}{solamente el punto \( X \)}{}{}}
\End

\Data\ID{1103059607}
\Updated{2021-01-15 09:01:22}
\Level{B}
\SubArea{Lines and planes: intersecting, perpendicular, parallel}
\IMG{1103059607.pdf}
\LANGen{
\Question{
Let \( ABCDV \) be a rectangle based-pyramid, where \( V \) is its apex, and let \( XY \) be a line where:
\begin{align*}
X&\text{ lays on a ray }BA\text{ and }|BA|=|AX|,\\
Y&\text{ lays on the height }SV\text{ and }|SY|=|YV|,\\
S&\text{ is the centre of the pyramid base}
\end{align*}
(see the picture). The points of intersection of the line \( XY \) with the surface of the pyramid lay:}
{on the pyramid faces \( ADV \) and \( BCV \)}{on the pyramid faces \( DCV \) and \( ABV \)}{on the pyramid face \( ADV \) and its edge \( CV \)}{on the pyramid edges \( AV \) and \( CV \)}{}{}}
\LANGpl{
\Question{
Dany jest ostrosłup \( ABCDV \), którego postawą jest prostokąt, gdzie punkt \( V \) to jego wierzchołek, prosta  \( XY \) to prosta, gdzie:
\begin{align*}
X&\text{ leży na półprostej }BA\text{ i }|BA|=|AX|,\\
Y&\text{ leży na wysokości }SV\text{ i }|SY|=|YV|,\\
S&\text{ to środek podstawy ostrosłupa}
\end{align*}
(spójrz na rysunek). Punkty przecięcia prostej \( XY \) z powierzchnią ostrosłupa leżą:}
{na ścianach \( ADV \) i \( BCV \)}{na ścianach \( DCV \) i \( ABV \)}{na ścianach \( ADV \) oraz krawędzi \( CV \)}{na krawędziach \( AV \) i \( CV \)}{}{}}
\LANGcs{
\Question{
Je dán pravidelný čtyřboký jehlan \( ABCDV \), kde \( V \) je hlavní vrchol jehlanu. Přímka \( XY \) je určena takto:
\begin{align*}
X&\text{ leží na polopřímce }BA\text{ a }|BA|=|AX|,\\
Y&\text{ leží na tělesové výšce jehlanu }SV\text{ a }|SY|=|YV|,\\
S&\text{ je střed podstavy jehlanu}
\end{align*}
(viz obrázek). Průsečíky přímky \( XY \) s povrchem jehlanu leží:}
{ve stěnách jehlanu \( ADV \) a \( BCV \)}{ve stěnách jehlanu \( DCV \) a \( ABV \)}{ve stěně jehlanu \( ADV \) a na hraně  \( CV \)}{na hranách jehlanu \( AV \) a \( CV \)}{}{}}
\LANGsk{
\Question{
Je daný pravidelný štvorboký ihlan \( ABCDV \), kde \( V \) je hlavný vrchol ihlanu. Priamka \( XY \) je určená takto:
\begin{align*}
X&\text{ leží na polpriamke }BA\text{ a }|BA|=|AX|,\\
Y&\text{ leží na telesovej výške ihlanu }SV\text{ a }|SY|=|YV|,\\
S&\text{ je stred podstavy ihlanu}
\end{align*}
(viď obrázok). Priesečníky priamky \( XY \) s povrchom ihlanu leží:}
{v stenách ihlanu \( ADV \) a \( BCV \)}{v stenách ihlanu \( DCV \) a \( ABV \)}{v stene ihlanu \( ADV \) a na hrane  \( CV \)}{na hranách ihlanu \( AV \) a \( CV \)}{}{}}
\LANGes{
\Question{
Sea un pirámide regular cuadrilátero \( ABCDV \), donde  \( V \) es el vértice del pirámide. La recta \( XY \) se define:
\begin{align*}
X&\text{ pertenece a la semirecta }BA\text{ y }|BA|=|AX|,\\
Y&\text{ pertenece a la altura del pirámide }SV\text{ y }|SY|=|YV|,\\
S&\text{ es el centro de la base del pirámide}
\end{align*}
(vea el dibujo). Las intersecciones de la recta \( XY \) con la superficie del pirámide pertenecen a:}
{los lados \( ADV \) y \( BCV \)}{los lados \( DCV \) y \( ABV \)}{el lado \( ADV \) y al borde  \( CV \)}{los bordes \( AV \) y \( CV \)}{}{}}
\End
